\documentclass[12pt,a4paper,twoside]{report}

% ======================
% Packages
% ======================
\usepackage[utf8]{inputenc}
\usepackage[T1]{fontenc}
\usepackage[french]{babel}
\usepackage{times}
\usepackage{geometry}
\geometry{margin=2.5cm}
\usepackage{setspace}
\setstretch{1.15}

\usepackage{graphicx}
\usepackage{caption}
\usepackage{float}
\usepackage{fancyhdr}
\usepackage{titlesec}
\usepackage{ragged2e}
\usepackage{tcolorbox}
\usepackage{hyperref}
\usepackage{lastpage}
\usepackage{xcolor}
\usepackage{booktabs}
\usepackage{array}
\usepackage{longtable}
\usepackage{enumitem}
\usepackage{amsmath}
\usepackage{listings}
\usepackage{tikz}
\usetikzlibrary{shapes.geometric, arrows, arrows.meta, positioning, calc,
  backgrounds, fit, decorations.pathreplacing, shadings}

% ======================
% Couleurs PwC
% ======================
\definecolor{pwcorange}{RGB}{208,74,2}
\definecolor{pwcblack}{RGB}{29,29,27}
\definecolor{pwcgray}{RGB}{100,100,100}
\definecolor{crisppink}{RGB}{220,50,120}
\definecolor{crispamber}{RGB}{255,165,0}

% ======================
% Pagination
% ======================
\pagestyle{fancy}
\fancyhf{}
\fancyhead[LE,RO]{\nouppercase{\leftmark}}
\fancyfoot[R]{\thepage/\pageref{LastPage}}
\renewcommand{\headrulewidth}{0.4pt}

% ======================
% Paragraphes
% ======================
\setlength{\parindent}{0.5cm}
\justifying

% ======================
% Titres
% ======================
\titleformat{\chapter}[hang]{\bfseries\fontsize{16}{18}\selectfont}{\thechapter}{0.2cm}{}
\titlespacing*{\chapter}{0pt}{12pt}{12pt}
\titleformat{\section}[hang]{\bfseries\fontsize{14}{16}\selectfont}{\thesection}{0.2cm}{}
\titlespacing*{\section}{0pt}{10pt}{6pt}
\titleformat{\subsection}[hang]{\bfseries\fontsize{12}{14}\selectfont}{\thesubsection}{0.2cm}{}
\titlespacing*{\subsection}{0pt}{8pt}{4pt}

% ======================
% Encadré définition
% ======================
\newtcolorbox{definitionbox}{
  colframe=black,
  colback=gray!5,
  boxrule=0.6pt,
  arc=2pt
}

% ======================
% Listings
% ======================
\lstset{
  basicstyle=\ttfamily\small,
  breaklines=true,
  frame=single,
  captionpos=b,
  numbers=left,
  numberstyle=\tiny,
  xleftmargin=1.5em
}

% ======================
% Métadonnées PDF
% ======================
\hypersetup{
  colorlinks=true,
  linkcolor=black,
  urlcolor=blue,
  citecolor=black,
  pdftitle={Détection intelligente d'anomalies financières},
  pdfauthor={Yousra Chaieb},
  pdfsubject={Projet de Fin d'Études --- Méthodologie CRISP-DM}
}

\begin{document}

\pagenumbering{gobble}

% =====================================================
% PAGE DE GARDE
% =====================================================
\begin{titlepage}
\centering

\vspace{1cm}

{\large \textbf{École Supérieure Privée d'Ingénierie et de Technologies --- ESPRIT}}\\[0.3cm]
{\large Département Informatique --- Filière Data Science \& IA}

\vspace{1.5cm}

\rule{\linewidth}{0.4pt}\\[0.4cm]
{\fontsize{24}{28}\selectfont \textbf{Détection intelligente\\[0.2cm]
d'anomalies financières}}\\[0.4cm]
\rule{\linewidth}{0.4pt}

\vspace{0.6cm}

{\large \textit{Analyse de transactions, modélisation par apprentissage automatique\\
et génération d'explications par modèle de langage}}

\vspace{1.5cm}

{\large Rapport de Projet de Fin d'Études}\\[0.2cm]
{\normalsize Réalisé dans le cadre du cycle ingénieur ---
Spécialité Intelligence Artificielle}

\vspace{1.5cm}

\begin{tabular}{rl}
  \textbf{Réalisé par :}           & Yousra Chaieb \\[0.2cm]
  \textbf{Encadrant académique :}  & Pr. Mohamed Islem Samaali \\[0.2cm]
  \textbf{Encadrants entreprise :} & Mr. Slim Zribi \& Mme. Hajer \\[0.2cm]
\end{tabular}

\vfill

{\large Année universitaire : 2025/2026}

\end{titlepage}

\clearpage
\pagenumbering{roman}

% =====================================================
% DÉDICACES
% =====================================================
\chapter*{Dédicaces}
\addcontentsline{toc}{chapter}{Dédicaces}

\vspace{3cm}

\begin{center}
\itshape\large

Je dédie ce modeste travail\\[0.6cm]

À mon père,\\
source inépuisable de force et de sagesse,\\
dont le soutien indéfectible m'a porté tout au long de ce chemin.\\[0.6cm]

À ma mère,\\
pour son amour inconditionnel, sa tendresse et ses prières\\
qui ont illuminé chacune de mes journées.\\[0.6cm]

À ma chère sœur \textbf{Eya},\\
pour sa complicité, son encouragement constant\\
et les moments de joie partagés.\\[0.6cm]

À mes amis,\\
pour leur présence, leur bonne humeur et leur précieux soutien\\
tout au long de cette aventure.\\[0.6cm]

À toutes les personnes qui m'ont encouragé et aidé\\
durant mon stage de fin d'études,\\
et qui ont contribué, de près ou de loin,\\
à la réussite de ce projet.

\end{center}

\newpage

% =====================================================
% REMERCIEMENTS
% =====================================================
\chapter*{Remerciements}
\addcontentsline{toc}{chapter}{Remerciements}

Après avoir rendu grâce à \textbf{Dieu}, le Tout-Puissant et le Miséricordieux, qui m'a
accordé la volonté, la patience et la santé nécessaires pour mener à bien ce projet de
fin d'études.

\medskip

Je tiens à exprimer ma profonde gratitude à mon encadrant académique,
\textbf{Mr. Mohamed Islem Samaali}, pour ses conseils avisés, sa disponibilité permanente
et ses orientations scientifiques qui ont été déterminantes pour structurer ma démarche
et enrichir mes réflexions tout au long de ce travail.

\medskip

Mes sincères remerciements s'adressent également à mes encadrants en entreprise,
\textbf{Mr. Slim} et \textbf{Madame Hajer}, pour leur accueil chaleureux, la confiance
qu'ils m'ont accordée et les précieuses orientations métier qui m'ont permis de mieux
appréhender les enjeux réels de la détection d'anomalies financières.

\medskip

Je remercie tout particulièrement \textbf{Madame Emna Zanni}, responsable du département
Audit IT, pour son soutien, sa bienveillance et les ressources qu'elle a mises à ma
disposition, contribuant ainsi de manière significative à la réussite de ce projet.

\medskip

Enfin, je remercie chaleureusement mes \textbf{camarades de stage} pour les échanges
enrichissants, la bonne ambiance et l'entraide qui ont rendu cette expérience encore plus
formatrice et agréable.

\newpage

% =====================================================
% RÉSUMÉ
% =====================================================
\chapter*{Résumé}
\addcontentsline{toc}{chapter}{Résumé}

Ce projet de fin d'études porte sur la détection intelligente d'anomalies financières à
partir de données transactionnelles réelles collectées auprès d'un client dans le cadre
d'une mission d'audit réalisée au sein de \textbf{PwC Tunisie}. Conduit selon la
méthodologie \textbf{CRISP-DM} (\textit{Cross-Industry Standard Process for Data Mining}),
il vise à construire une chaîne complète et reproductible allant de la compréhension métier
et des données jusqu'au déploiement d'un pipeline de détection opérationnel.

Le corpus utilisé comprend environ \textbf{200~000 transactions financières}, dont
258 opérations frauduleuses (ratio 1:774). La sévérité du déséquilibre des classes
représente le défi central du projet et justifie l'emploi de stratégies spécifiques
telles que la pondération des classes, le rééchantillonnage SMOTE et une évaluation
fondée sur le rappel, la précision et le F1-score plutôt que sur la simple exactitude.

Trois approches de modélisation sont comparées : (1) une \textbf{régression logistique}
servant de baseline linéaire, (2) une \textbf{forêt aléatoire} constituant la référence
non-linéaire supervisée, et (3) un \textbf{autoencoder} profond non supervisé entraîné
exclusivement sur les transactions légitimes.

Afin d'améliorer l'interprétabilité, une double couche explicative a été mise en place :
les méthodes \textbf{SHAP} et \textbf{LIME} pour identifier les variables les plus
contributives, et une intégration avec \textbf{Ollama} pour générer des explications
en langage naturel accessibles aux analystes métier.

\vspace{0.5cm}
\textbf{Mots-clés :} détection d'anomalies, fraude financière, CRISP-DM, autoencoder,
forêt aléatoire, SMOTE, SHAP, LIME, Ollama, interprétabilité, PwC Tunisie.

\newpage

% =====================================================
% ABSTRACT
% =====================================================
\chapter*{Abstract}
\addcontentsline{toc}{chapter}{Abstract}

This final-year project focuses on intelligent anomaly detection in financial transactions
using a dataset of real transactional records collected from a client engagement conducted
within \textbf{PwC Tunisia}. Following the \textbf{CRISP-DM} methodology, the goal is to
build a complete and reproducible pipeline covering business understanding, data exploration,
preparation, modeling, evaluation, and deployment.

The dataset comprises approximately \textbf{200,000 financial transactions}, of which only
258 are fraudulent (ratio 1:774). Three modeling approaches are benchmarked:
(1) \textbf{logistic regression} as a linear baseline, (2) \textbf{random forest} as a
robust supervised reference, and (3) a deep \textbf{autoencoder} trained exclusively on
legitimate transactions. To enhance interpretability, a two-layer explainability stack
was built: \textbf{SHAP} and \textbf{LIME} methods, and an \textbf{Ollama} integration
to generate natural-language explanations accessible to business analysts.

\vspace{0.5cm}
\textbf{Keywords:} anomaly detection, financial fraud, CRISP-DM, autoencoder, random
forest, SMOTE, SHAP, LIME, Ollama, explainability, PwC Tunisia.

\newpage

% =====================================================
% TABLE DES MATIÈRES
% =====================================================
\tableofcontents
\newpage

% =====================================================
% LISTE DES FIGURES
% =====================================================
\listoffigures
\addcontentsline{toc}{chapter}{Liste des figures}
\newpage

% =====================================================
% LISTE DES TABLEAUX
% =====================================================
\listoftables
\addcontentsline{toc}{chapter}{Liste des tableaux}
\newpage

% =====================================================
% LISTE DES ABRÉVIATIONS
% =====================================================
\chapter*{Liste des abréviations}
\addcontentsline{toc}{chapter}{Liste des abréviations}

\begin{tabular}{ll}
  \textbf{AE}       & AutoEncoder \\
  \textbf{AUC}      & Area Under the Curve \\
  \textbf{BN}       & Batch Normalization \\
  \textbf{CRISP-DM} & Cross-Industry Standard Process for Data Mining \\
  \textbf{EDA}      & Exploratory Data Analysis (Analyse Exploratoire des Données) \\
  \textbf{ERP}      & Enterprise Resource Planning (Progiciel de Gestion Intégré) \\
  \textbf{F1}       & F1-Score \\
  \textbf{FN}       & Faux Négatifs \\
  \textbf{FP}       & Faux Positifs \\
  \textbf{LLM}      & Large Language Model (Grand Modèle de Langage) \\
  \textbf{LIME}     & Local Interpretable Model-agnostic Explanations \\
  \textbf{LR}       & Logistic Regression (Régression Logistique) \\
  \textbf{ML}       & Machine Learning (Apprentissage Automatique) \\
  \textbf{MSE}      & Mean Squared Error (Erreur Quadratique Moyenne) \\
  \textbf{ONG}      & Organisation Non Gouvernementale \\
  \textbf{PR-AUC}   & Area Under the Precision-Recall Curve \\
  \textbf{RF}       & Random Forest (Forêt Aléatoire) \\
  \textbf{RCM}      & Risk \& Control Matrix (Matrice de Risques et de Contrôles) \\
  \textbf{ROC}      & Receiver Operating Characteristic \\
  \textbf{SHAP}     & SHapley Additive exPlanations \\
  \textbf{SMOTE}    & Synthetic Minority Over-sampling Technique \\
  \textbf{TP}       & Vrais Positifs \\
  \textbf{TN}       & Vrais Négatifs \\
  \textbf{VAE}      & Variational AutoEncoder \\
  \textbf{XAI}      & Explainable Artificial Intelligence \\
\end{tabular}

\newpage
\pagenumbering{arabic}

% =====================================================
% INTRODUCTION GÉNÉRALE
% =====================================================
\chapter*{Introduction générale}
\addcontentsline{toc}{chapter}{Introduction générale}

Au cours de mon stage au sein du département Risk Assurance Services de PwC Tunisie,
j'ai été confrontée à une réalité concrète : dans le cadre d'une mission d'audit IT,
l'équipe dispose de plusieurs centaines de milliers de lignes transactionnelles à analyser,
et seule une poignée d'entre elles dissimule une anomalie. Repérer ces cas dans une masse
de données aussi volumineuse, sans passer à côté d'une fraude ni submerger les auditeurs
de fausses alertes --- voilà le défi auquel répond ce projet.

Ce travail propose une \textbf{chaîne de traitement complète} pour la détection intelligente
d'anomalies financières. Développé selon la méthodologie \textbf{CRISP-DM}, il couvre
l'intégralité du cycle : de la compréhension des besoins métier à la construction d'un
pipeline de déploiement reproductible, en passant par l'analyse exploratoire des données,
leur préparation, l'entraînement de plusieurs modèles et leur évaluation rigoureuse.

Deux défis techniques ont structuré l'ensemble de la démarche : (1) la \textbf{gestion
du déséquilibre extrême} des classes --- moins de 0,13\,\% des transactions sont
frauduleuses --- et (2) l'\textbf{interprétabilité} des décisions, rendue possible grâce
aux méthodes SHAP, LIME et à la génération d'explications textuelles par le modèle de
langage Ollama.

\paragraph{Organisation du rapport :} Ce document est structuré en sept chapitres :
\begin{itemize}
  \item Le \textbf{chapitre~1} présente l'organisme d'accueil PwC Tunisie, le cadre du
    stage et les objectifs du projet.
  \item Le \textbf{chapitre~2} établit la compréhension métier du problème, la
    problématique et les critères de succès (phase~1 CRISP-DM).
  \item Le \textbf{chapitre~3} présente le jeu de données transactionnel et l'analyse
    exploratoire réalisée (phase~2 CRISP-DM).
  \item Le \textbf{chapitre~4} détaille les opérations de préparation des données :
    nettoyage, ingénierie de variables et gestion du déséquilibre (phase~3 CRISP-DM).
  \item Le \textbf{chapitre~5} décrit les modèles construits : régression logistique,
    forêt aléatoire et autoencoder (phase~4 CRISP-DM).
  \item Le \textbf{chapitre~6} présente l'évaluation comparative des modèles et les
    mécanismes d'explicabilité (phase~5 CRISP-DM).
  \item Le \textbf{chapitre~7} aborde le déploiement, l'architecture technique et la
    reproductibilité du pipeline (phase~6 CRISP-DM).
\end{itemize}

% =====================================================
% CHAPITRES (inclus depuis chapitres/)
% =====================================================
% =====================================================
% CHAPITRE 1 : PRÉSENTATION DE L'ORGANISME D'ACCUEIL
% =====================================================
\chapter{Présentation de l'organisme d'accueil}

% ---------------------------------------------------
\section*{Introduction}
\addcontentsline{toc}{section}{Introduction}
% ---------------------------------------------------

Ce premier chapitre pose le cadre institutionnel et professionnel dans lequel s'inscrit
ce projet de fin d'études. Avant de décrire les choix techniques et les résultats obtenus,
il est indispensable de comprendre l'environnement au sein duquel la problématique a émergé.
Nous présenterons successivement le réseau PwC à l'échelle internationale, son implantation
tunisienne, puis le département Risk Assurance Services et sa cellule Audit IT --- contexte
direct de production des données utilisées dans ce travail. Nous conclurons par la description
du cadre du stage et des missions confiées.

% ---------------------------------------------------
\section{PwC à l'échelle internationale}
% ---------------------------------------------------

\subsection{Historique et positionnement}

PwC est un réseau mondial de cabinets de services professionnels couvrant l'audit, le conseil,
la fiscalité et les transactions. Sa clientèle est très diversifiée : grands groupes
internationaux cotés, PME en pleine croissance, institutions publiques et organisations à but
non lucratif. Ce qui caractérise le réseau, c'est sa capacité à articuler une expertise
internationale uniformisée avec une connaissance fine des réalités locales --- une posture
particulièrement pertinente sur des marchés comme la Tunisie, où les spécificités
réglementaires et économiques sont fortes.

PwC est né en \textbf{1998} de la fusion de deux cabinets alors centenaires :
\textit{Price Waterhouse}, fondé à Londres en 1849, et \textit{Coopers \& Lybrand},
dont les origines remontent à 1854. L'entité fusionnée, d'abord connue sous le nom
PricewaterhouseCoopers, a adopté en septembre 2010 son identité actuelle abrégée « PwC ».
Aujourd'hui présent dans plus de \textbf{187 pays} avec \textbf{771 bureaux} et
\textbf{364~000 collaborateurs}, le réseau se positionne comme le premier cabinet mondial
d'audit et de conseil, devant Deloitte, EY et KPMG.

\subsection{PwC dans le Big Four}

PwC s'inscrit dans le cercle restreint des quatre grands cabinets mondiaux d'audit et de
conseil --- le \og Big Four \fg{} ---, aux côtés de Deloitte, Ernst \& Young et KPMG.
Chacun de ces réseaux repose sur le même modèle organisationnel : des entités nationales
juridiquement indépendantes, fédérées sous une marque commune et soumises à des standards
partagés en matière de qualité et d'éthique professionnelle.

\begin{table}[H]
\centering
\caption{Fiche signalétique de PwC International}
\label{tab:fiche_pwc_intl}
\begin{tabular}{ll}
\toprule
\textbf{Critère}              & \textbf{Information} \\
\midrule
Dénomination                  & PricewaterhouseCoopers (PwC) \\
Date de création              & 1998 (fusion Price Waterhouse \& Coopers \& Lybrand) \\
Siège social                  & Londres, Royaume-Uni \\
Secteur d'activité            & Audit, Conseil, Fiscalité, Transactions \\
Nombre de pays                & Plus de 187 \\
Nombre de bureaux             & Plus de 771 \\
Effectif mondial              & Plus de 364~000 collaborateurs \\
Appartenance                  & Big Four de l'audit \\
\bottomrule
\end{tabular}
\end{table}

\subsection{Statut juridique}

Sur le plan juridique, PwC s'appuie sur \textit{PricewaterhouseCoopers International
Limited} (PwCIL), une société de droit anglais qui joue un rôle de coordination
stratégique sans être elle-même prestataire de services. Chaque cabinet national est
une entité indépendante, ce qui lui permet de satisfaire aux exigences des pays imposant
un actionnariat local pour les cabinets d'audit. En contrepartie de l'usage de la marque
PwC, ces entités s'engagent à respecter les standards de qualité, de gestion du risque
et d'éthique professionnelle définis par PwCIL.

\subsection{Stratégie mondiale : « The New Equation »}

Pour anticiper les évolutions du marché --- notamment l'essor de l'intelligence
artificielle et les nouvelles exigences réglementaires ---, PwC a structuré sa stratégie
globale autour du programme \textbf{« The New Equation »}, articulé autour de deux axes :

\begin{itemize}
  \item \textbf{Building Trust} : répondre aux attentes en matière de transparence en
    combinant expertise d'audit, de fiscalité et de cybersécurité, enrichies par
    l'intelligence artificielle.
  \item \textbf{Sustained Outcomes} : mobiliser l'ensemble des expertises du réseau pour
    accompagner ses clients vers des résultats durables, couvrant stratégie, services
    numériques, fiscalité et conformité réglementaire.
\end{itemize}

% ---------------------------------------------------
\section{PwC Tunisie}
% ---------------------------------------------------

\subsection{Présentation générale}

Fondée en \textbf{1983}, PwC Tunisie s'impose comme l'une des premières sociétés
tunisiennes dans le domaine de l'audit et du conseil, forte de plus de quarante années
d'expérience au service du tissu économique local. Elle propose une gamme étendue de
services adaptés au marché tunisien, couvrant aussi bien les besoins des institutions
publiques que ceux des entreprises privées, des ONG et des bailleurs de fonds
internationaux.

\begin{table}[H]
\centering
\caption{Fiche signalétique de PwC Tunisie}
\label{tab:fiche_pwc_tunisie}
\begin{tabular}{ll}
\toprule
\textbf{Critère}              & \textbf{Information} \\
\midrule
Dénomination                  & PricewaterhouseCoopers Tunisie \\
Date de création              & 1983 \\
Siège social                  & Tunis, Tunisie \\
Secteur d'activité            & Audit, Conseil, Fiscalité, Transactions \\
Domaines d'intervention       & Secteur public, Privé, ONG, Bailleurs de fonds \\
Appartenance                  & Réseau PwC International \\
\bottomrule
\end{tabular}
\end{table}

\subsection{Offres de services de PwC Tunisie}

PwC Tunisie s'organise en quatre lignes de services complémentaires :

\paragraph{Audit.}
Évaluation du contrôle interne, conception d'organigrammes, développement de manuels de
procédures, construction de matrices de risques et de contrôles (Risk \& Control
Matrix --- RCM).

\paragraph{Transactions (Deals).}
Due diligence acheteur et vendeur, études de faisabilité, élaboration de business plans,
évaluation d'entreprises et gestion des processus d'ouverture de capital.

\paragraph{Services Juridiques et Fiscaux.}
Conseil fiscal, tax due diligence, audit fiscal, tax compliance et assistance lors des
contrôles fiscaux.

\paragraph{Consulting.}
Définition de stratégies de développement, refonte commerciale, amélioration de
l'efficacité opérationnelle et conseil en marketing.

\subsection{Références clients}

PwC Tunisie opère dans tous les domaines d'activité, avec une expertise forte dans les
services financiers, le secteur public, l'énergie, les projets d'infrastructure et les
télécommunications.

% ---------------------------------------------------
\section{Le département Risk Assurance Services}
% ---------------------------------------------------

\subsection{Présentation générale}

Le département \textbf{Risk Assurance Services} (anciennement Conseil Audit Formation) est
le pôle au sein duquel j'ai réalisé ce stage. Sa mission centrale : aider les organisations
clientes à identifier et cartographier leurs risques opérationnels et systémiques, puis à
mettre en place les dispositifs de contrôle appropriés. Concrètement, les missions couvrent
un spectre large --- de l'évaluation du contrôle interne à la revue des systèmes
d'information, en passant par l'assistance à la conception de procédures et la formation des
équipes audit.

\subsection{Offres de prestations}

Le département apporte une valeur ajoutée à travers trois formes d'intervention
complémentaires :

\paragraph{Outsourcing des travaux.}
Cette prestation s'adresse aux organisations ne disposant pas d'une fonction d'audit interne
constituée. PwC prend alors en charge l'intégralité des travaux d'audit, assurant une
couverture complète de l'ensemble des missions inscrites au plan d'audit interne.

\paragraph{Co-sourcing des travaux.}
L'équipe d'audit interne du client fait appel à l'expertise de PwC pour renforcer ses
capacités. Cette modalité collaborative favorise le transfert de connaissances et permet à
l'équipe interne de monter progressivement en compétences.

\paragraph{Services de conseil.}
PwC accompagne les entreprises dans l'évaluation de leurs opérations et dans la structuration
ou l'amélioration de leur fonction d'audit interne, afin de la rendre plus performante et
mieux alignée sur leurs enjeux stratégiques.

\subsection{La cellule Audit IT}

Au cœur du département Risk Assurance Services, la cellule \textbf{Audit IT} se consacre
à l'examen des systèmes d'information qui alimentent la comptabilité et les états financiers
des clients : ERP, systèmes de gestion de trésorerie, plateformes de paiement. Son rôle
est de vérifier que ces environnements produisent des données intègres, fiables et conformes
aux normes d'audit --- ce qui conditionne directement la validité des assertions financières
sur lesquelles repose l'opinion du commissaire aux comptes.

C'est précisément dans ce cadre opérationnel que les données de ce projet ont été collectées :
un corpus de transactions financières extraites du système d'information d'un client de PwC
Tunisie, dans le cadre d'une mission d'audit IT active. Pour des raisons de confidentialité,
toutes les informations permettant d'identifier les parties ont été anonymisées.

% ---------------------------------------------------
\section{Cadre du stage et projet de fin d'études}
% ---------------------------------------------------

\subsection{Objectifs du stage}

Ce stage de fin d'études, réalisé au sein du département Risk Assurance Services de PwC
Tunisie, a pour objectif principal de concevoir et de déployer un \textbf{système de
détection intelligente d'anomalies financières} fondé sur des techniques avancées
d'apprentissage automatique. Il s'inscrit dans une démarche d'innovation visant à renforcer
les capacités analytiques de l'équipe d'audit, en automatisant partiellement l'identification
des transactions suspectes et en produisant des explications compréhensibles par les
analystes métier.

\subsection{Missions confiées}

Les missions confiées dans le cadre de ce stage couvrent l'intégralité du cycle CRISP-DM :

\begin{itemize}
  \item Compréhension du contexte métier de la détection de fraude financière en
    environnement d'audit ;
  \item Analyse exploratoire du corpus de transactions réelles
    ($\sim$200~000 opérations, ratio fraude 1:774) ;
  \item Conception et implémentation d'un pipeline de préparation des données
    (nettoyage, ingénierie de variables, gestion du déséquilibre via SMOTE) ;
  \item Développement et comparaison de trois approches de modélisation : régression
    logistique, forêt aléatoire et autoencoder non supervisé ;
  \item Mise en place d'une couche d'explicabilité double : méthodes SHAP et LIME, et
    génération d'explications en langage naturel via \textbf{Ollama} ;
  \item Déploiement d'un pipeline reproductible via le script d'orchestration
    \texttt{run\_all.py}, avec couverture de tests unitaires.
\end{itemize}

\subsection{Présentation du projet}

Le projet vise à construire une \textbf{chaîne de traitement complète et reproductible} pour
la détection de transactions frauduleuses. Conduit selon la méthodologie \textbf{CRISP-DM},
il répond à la problématique suivante : \textit{comment détecter efficacement les fraudes
dans un corpus fortement déséquilibré, tout en produisant des explications compréhensibles
pour les auditeurs ?}

La solution développée repose sur trois contributions complémentaires :

\begin{itemize}
  \item \textbf{Modélisation multi-approches} : comparaison d'un modèle linéaire, d'un
    modèle ensembliste supervisé et d'un autoencoder profond non supervisé, évalués sur
    des métriques adaptées au déséquilibre (F1-score, PR-AUC, rappel) ;
  \item \textbf{Gestion du déséquilibre} : stratégies de pondération des classes et de
    rééchantillonnage SMOTE pour traiter le ratio fraude/légitime de 1:774 ;
  \item \textbf{Explicabilité des décisions} : identification des variables les plus
    contributives via SHAP et LIME, et génération d'explications en langage naturel à
    destination des analystes métier via le LLM Ollama.
\end{itemize}

% ---------------------------------------------------
\section*{Conclusion}
\addcontentsline{toc}{section}{Conclusion}
% ---------------------------------------------------

Ce premier chapitre a établi le cadre institutionnel et organisationnel dans lequel s'inscrit
ce projet de fin d'études. À travers la présentation de PwC en tant que réseau international
de référence, puis de PwC Tunisie et de son département Risk Assurance Services, nous avons
mis en lumière l'environnement professionnel au sein duquel ce stage a été effectué, ainsi
que le rôle stratégique de la cellule Audit IT. La problématique identifiée --- détecter
automatiquement les transactions frauduleuses dans un corpus réel fortement déséquilibré,
tout en garantissant l'interprétabilité des résultats --- et les missions confiées ont posé
les bases nécessaires à la compréhension des enjeux du projet.

Le chapitre suivant sera consacré à la compréhension métier, première phase de la
méthodologie CRISP-DM, dans laquelle nous formaliserons la problématique, les objectifs
et les critères de succès du projet.

% =====================================================
% CHAPITRE 2 : COMPRÉHENSION MÉTIER (CRISP-DM Phase 1)
% =====================================================
\chapter{Compréhension métier}

% ---------------------------------------------------
\section*{Introduction}
\addcontentsline{toc}{section}{Introduction}
% ---------------------------------------------------

Avant d'écrire la moindre ligne de code, il faut comprendre le problème que l'on cherche
à résoudre --- et surtout, le métier derrière les données. C'est l'essence de la première
phase CRISP-DM. Ce chapitre pose donc le cadre : contexte de la fraude financière telle
qu'on la rencontre en mission d'audit, formalisation de la problématique, objectifs
opérationnels et critères de succès mesurables qui guideront toutes les décisions techniques
ultérieures. Cette étape est non négociable dans un contexte d'audit IT : une définition
floue des objectifs aurait conduit à des choix algorithmiques déconnectés des besoins réels
des auditeurs.

% ---------------------------------------------------
\section{Contexte de la détection de fraude financière}
% ---------------------------------------------------

\subsection{Enjeux économiques et opérationnels}

La fraude financière représente une menace persistante et coûteuse pour les institutions.
Si les chiffres globaux donnent le vertige --- le rapport Nilson (2022) évalue à plusieurs
dizaines de milliards de dollars les pertes annuelles liées à la fraude par carte bancaire
---, c'est à l'échelle d'une mission d'audit que l'impact se fait le plus ressentir :
transactions non reconnues, coûts de remboursement, investigations manuelles chronophages
et responsabilité engagée des auditeurs. À cela s'ajoutent des contraintes réglementaires
croissantes, comme la directive européenne PSD2 qui impose des mécanismes de surveillance
renforcés des transactions.

Pour une équipe comme celle du Risk Assurance Services de PwC Tunisie, l'enjeu est double :
détecter les anomalies avant qu'elles ne se propagent, et le faire de manière traçable et
défendable dans un rapport d'audit.

\subsection{Limites des systèmes basés sur des règles}

Historiquement, la détection de fraude repose sur des règles métier définies manuellement
par des experts : seuils de montant, listes noires de comptes suspects, patterns connus de
détournement. Ces systèmes ont le mérite d'être explicables et rapides à déployer. Mais
lors de l'analyse du corpus de ce stage, le constat a été frappant : le système naïf
\texttt{isFlaggedFraud} présent dans les données n'identifie qu'un seul cas de fraude sur
258, soit un rappel de \textbf{0,39\,\%}. Trois limites fondamentales expliquent cet
échec :

\begin{itemize}
  \item \textbf{Rigidité} : les règles codifient les comportements connus, mais les
    fraudeurs adaptent continuellement leurs méthodes pour les contourner.
  \item \textbf{Maintenance coûteuse} : chaque nouvelle forme de fraude nécessite une
    intervention manuelle d'un expert pour mettre à jour les règles.
  \item \textbf{Faux positifs élevés} : les seuils statiques génèrent de nombreuses
    alertes injustifiées, surchargeant les équipes d'investigation et dégradant la
    confiance dans le système.
\end{itemize}

L'apprentissage automatique offre une alternative complémentaire : plutôt que d'encoder les
patterns de fraude à la main, il les extrait automatiquement des données historiques, y
compris des interactions non linéaires entre variables que l'œil humain ne perçoit pas.

% ---------------------------------------------------
\section{Présentation et formalisation du projet}
% ---------------------------------------------------

\subsection{Contexte du projet de fin d'études}

Ce projet a été réalisé dans le cadre d'un stage de fin d'études axé sur la mise en œuvre
d'un système de détection d'anomalies financières exploitant des techniques avancées
d'apprentissage automatique. Le travail couvre l'intégralité du cycle de vie d'un projet
data science, de la définition du problème au déploiement d'un pipeline reproductible.

\subsection{Problématique}

\begin{definitionbox}
\textbf{Problématique :} Comment concevoir un pipeline reproductible, fondé sur la
méthodologie CRISP-DM, capable de détecter efficacement les transactions financières
frauduleuses dans un contexte de fort déséquilibre des classes, tout en produisant des
explications compréhensibles pour les analystes métier ?
\end{definitionbox}

Cette problématique est structurée autour de deux défis techniques majeurs : (1) la
\textbf{gestion du déséquilibre extrême} des classes --- moins de 0,13\,\% des transactions
sont frauduleuses, un ratio si faible qu'un modèle naïf prédisant systématiquement
« légitime » afficherait 99,87\,\% de précision tout en ne détectant aucune fraude ---
et (2) l'\textbf{interprétabilité} des décisions, rendue possible grâce aux méthodes SHAP,
LIME et à la génération d'explications textuelles par le modèle de langage Ollama.

\subsection{Objectifs du projet}

Les objectifs opérationnels définis en début de projet sont les suivants :

\begin{enumerate}
  \item Construire une \textbf{chaîne de traitement complète et reproductible} couvrant
    toutes les phases CRISP-DM.
  \item Réaliser une \textbf{analyse exploratoire approfondie} des données transactionnelles.
  \item Concevoir et appliquer un \textbf{pipeline de préparation des données} robuste,
    préservant l'intégrité temporelle et évitant toute fuite d'information
    (\textit{data leakage}).
  \item \textbf{Entraîner et comparer} au moins trois approches de détection : baseline
    linéaire, modèle ensembliste supervisé et autoencoder non supervisé.
  \item Évaluer les modèles avec des \textbf{métriques adaptées au déséquilibre} : rappel,
    précision, F1-score, courbe PR et PR-AUC.
  \item Produire des \textbf{explications interprétables} des décisions (SHAP, LIME, Ollama).
  \item Livrer un \textbf{code modulaire, documenté et testé} facilitant la maintenance
    et l'extension.
\end{enumerate}

\subsection{Critères de succès}

\begin{table}[H]
\centering
\caption{Critères de succès du projet}
\label{tab:criteres_succes}
\begin{tabular}{lll}
\toprule
\textbf{Critère} & \textbf{Valeur cible} & \textbf{Justification} \\
\midrule
Recall (test)    & $> 0{,}70$ & Priorité à la détection des fraudes \\
F1-score (test)  & $> 0{,}70$ & Équilibre précision/rappel \\
PR-AUC (test)    & $> 0{,}75$ & Synthèse du tradeoff \\
Reproductibilité & 100\,\%   & Pipeline exécutable via \texttt{run\_all.py} \\
Couverture tests & $> 80\,\%$ & Robustesse du code \\
\bottomrule
\end{tabular}
\end{table}

% ---------------------------------------------------
\section{Cadre méthodologique : CRISP-DM}
% ---------------------------------------------------

\subsection{Présentation de la méthodologie}

\begin{definitionbox}
\textbf{CRISP-DM} (\textit{Cross-Industry Standard Process for Data Mining}) est un cadre
méthodologique ouvert proposé en 1996 par un consortium réunissant SPSS, Daimler-Benz,
NCR et OHRA. Son objectif : structurer de façon reproductible les projets de fouille de
données, quelle que soit l'industrie ou les outils utilisés. Il organise le travail en
\textbf{six phases interdépendantes} formant un cycle itératif, du cadrage métier jusqu'à
la mise en production.
\end{definitionbox}

Les six phases de CRISP-DM sont :

\begin{enumerate}
  \item \textbf{Business Understanding} --- Compréhension métier et définition des objectifs.
  \item \textbf{Data Understanding} --- Collecte, exploration et évaluation de la qualité
    des données.
  \item \textbf{Data Preparation} --- Sélection, nettoyage, construction et formatage des
    données.
  \item \textbf{Modeling} --- Sélection, construction et ajustement des modèles.
  \item \textbf{Evaluation} --- Évaluation des résultats au regard des objectifs métier.
  \item \textbf{Deployment} --- Mise en production et documentation.
\end{enumerate}

\subsection{Justification du choix de CRISP-DM}

Plusieurs méthodologies structurées existent pour conduire des projets de data science.
Le tableau ci-dessous les compare pour justifier le choix retenu.

\begin{table}[H]
\centering
\caption{Comparaison des principales méthodologies de data science}
\label{tab:methodo_compare}
\begin{tabular}{p{2.3cm}p{3.4cm}p{3.4cm}p{3.2cm}}
\toprule
\textbf{Méthodologie} & \textbf{Points forts} & \textbf{Limites}
  & \textbf{Adéquation} \\
\midrule
\textbf{CRISP-DM} &
  Agnostique outils ; cycle itératif ; couverture complète (métier $\to$ déploiement) &
  Moins adapté aux projets agiles purs &
  \textbf{Excellent} \\[0.3cm]
\textbf{SEMMA} (SAS) &
  Simple et concis &
  Lié à SAS ; absence de phase déploiement &
  \textbf{Insuffisant} \\[0.3cm]
\textbf{KDD} &
  Fondement académique solide &
  Peu flexible ; orienté recherche &
  \textbf{Partiel} \\[0.3cm]
\textbf{Agile DS} &
  Itérations courtes, DevOps &
  Suppose équipe multidisciplinaire et CI/CD &
  \textbf{Inadapté} \\
\bottomrule
\end{tabular}
\end{table}

Le choix de CRISP-DM pour ce projet n'était pas arbitraire. Plusieurs raisons concrètes
l'ont motivé :

\begin{itemize}
  \item \textbf{Orientation métier explicite} : avant de toucher aux données, CRISP-DM
    impose de formaliser les objectifs et les critères de succès --- étape non négociable
    en contexte d'audit IT.
  \item \textbf{Couverture complète du cycle} : contrairement à SEMMA ou KDD, CRISP-DM
    couvre aussi bien la compréhension des données que le déploiement du pipeline.
  \item \textbf{Caractère itératif} : la préparation des données et la modélisation ne
    sont pas séquentielles en pratique. Les allers-retours entre phases sont formalisés
    et légitimés par CRISP-DM.
  \item \textbf{Indépendance des outils} : CRISP-DM étant agnostique, il s'adapte sans
    difficulté à une stack Python / scikit-learn / PyTorch / Ollama.
  \item \textbf{Lisibilité pour les équipes PwC} : CRISP-DM étant largement adopté dans
    les cabinets de conseil, les encadrants entreprise ont pu suivre la progression du
    projet sans avoir à maîtriser les détails techniques.
\end{itemize}

\subsection{Application au projet}

L'ensemble du rapport est structuré selon ces six phases. Le caractère itératif de
CRISP-DM se manifeste notamment dans les allers-retours entre l'analyse exploratoire
(phase 2) et la préparation des données (phase 3), ainsi qu'entre la modélisation
(phase 4) et l'évaluation (phase 5) lors du réglage des seuils de détection.

% ---------------------------------------------------
\section*{Conclusion}
\addcontentsline{toc}{section}{Conclusion}
% ---------------------------------------------------

Ce chapitre a établi le contexte métier du projet, la problématique à résoudre et les
objectifs à atteindre selon le cadre CRISP-DM. Deux défis principaux guideront toutes
les décisions techniques : la gestion du déséquilibre extrême des classes (ratio 1:774)
et la nécessité de produire des décisions interprétables pour les auditeurs. Le choix
de CRISP-DM comme méthodologie directrice a été justifié par sa couverture complète du
cycle de vie, sa nature itérative et son indépendance vis-à-vis des outils. Le chapitre
suivant présente les données sur lesquelles repose l'intégralité du travail.

% =====================================================
% CHAPITRE 3 : COMPRÉHENSION DES DONNÉES (CRISP-DM Phase 2)
% =====================================================
\chapter{Compréhension des données}

% ---------------------------------------------------
\section*{Introduction}
\addcontentsline{toc}{section}{Introduction}
% ---------------------------------------------------

La phase de compréhension des données constitue le socle de tout projet de data science.
Elle précède et conditionne chaque décision de modélisation : un modèle ne peut être meilleur
que les données sur lesquelles il est entraîné. Avant d'appliquer le moindre algorithme,
il est indispensable de saisir la structure, les limites et les opportunités du corpus
disponible. Ce chapitre correspond au notebook \texttt{01\_data\_understanding.ipynb} du
projet et présente le jeu de données, son analyse exploratoire, l'étude de l'existant,
ainsi que l'état de l'art des approches de détection d'anomalies financières.

% ---------------------------------------------------
\section{Présentation du jeu de données}
% ---------------------------------------------------

\subsection{Origine et description}

Le jeu de données utilisé dans ce projet provient d'une \textbf{mission d'audit IT réalisée
au sein de PwC Tunisie}. Il regroupe des transactions financières issues du système
d'information d'un client. Pour des raisons de confidentialité, toutes les informations
permettant d'identifier les parties concernées ont été anonymisées.

\subsection{Caractéristiques générales}

Le corpus est constitué d'environ \textbf{200~000 transactions financières} enregistrées
sur une période couvrant 743 heures (soit environ 31 jours). Parmi ces transactions, seules
\textbf{258 sont frauduleuses}, soit un taux de fraude de \textbf{0,129\,\%}.
Ce déséquilibre extrême (ratio 1:774) constitue le défi central du projet.

\subsection{Description des variables}

\begin{table}[H]
\centering
\caption{Description des variables du jeu de données}
\label{tab:variables}
\begin{tabular}{lll}
\toprule
\textbf{Variable} & \textbf{Type} & \textbf{Description} \\
\midrule
\texttt{step}           & Entier    & Pas de temps (1 = 1 heure, max 743) \\
\texttt{type}           & Catégorie & Type d'opération (5 modalités) \\
\texttt{amount}         & Réel      & Montant de la transaction \\
\texttt{nameOrig}       & Chaîne    & Identifiant du compte émetteur \\
\texttt{oldbalanceOrg}  & Réel      & Solde émetteur avant transaction \\
\texttt{newbalanceOrig} & Réel      & Solde émetteur après transaction \\
\texttt{nameDest}       & Chaîne    & Identifiant du compte destinataire \\
\texttt{oldbalanceDest} & Réel      & Solde destinataire avant transaction \\
\texttt{newbalanceDest} & Réel      & Solde destinataire après transaction \\
\texttt{isFraud}        & Binaire   & \textbf{Variable cible} : 1 = fraude \\
\texttt{isFlaggedFraud} & Binaire   & Indicateur système naïf existant \\
\bottomrule
\end{tabular}
\end{table}

\subsection{Types de transactions}

Cinq types de transactions sont présents dans le jeu de données :

\begin{table}[H]
\centering
\caption{Distribution des types de transactions et taux de fraude}
\label{tab:types_transactions}
\begin{tabular}{lrr}
\toprule
\textbf{Type} & \textbf{Fréquence} & \textbf{Taux de fraude} \\
\midrule
CASH\_IN   & $\sim$35\,\% & 0,000\,\% \\
CASH\_OUT  & $\sim$35\,\% & 0,181\,\% \\
DEBIT      & $\sim$1\,\%  & 0,000\,\% \\
PAYMENT    & $\sim$22\,\% & 0,000\,\% \\
TRANSFER   & $\sim$7\,\%  & 0,775\,\% \\
\bottomrule
\end{tabular}
\end{table}

\textbf{Observation clé :} \textit{seuls les types TRANSFER et CASH\_OUT présentent des cas
de fraude}. Cette information sera directement exploitée lors de l'ingénierie de variables
(création de la feature \texttt{is\_transfer\_or\_cashout}).

% ---------------------------------------------------
\section{Analyse exploratoire des données (EDA)}
% ---------------------------------------------------

\subsection{Déséquilibre des classes}

Le déséquilibre est extrême : les transactions frauduleuses représentent \textbf{0,129\,\%}
du corpus. Un modèle prédisant systématiquement « légitime » atteindrait une exactitude de
99,87\,\% tout en ne détectant \textbf{aucune fraude}. C'est pourquoi l'exactitude
(\textit{accuracy}) est une métrique totalement inadaptée à ce problème, et que les
métriques de rappel, précision et F1-score ont été retenues comme indicateurs de succès.

\subsection{Distribution des montants}

La variable \texttt{amount} présente une asymétrie (\textit{skewness}) très élevée
(\textbf{30,80}), indiquant une distribution fortement concentrée sur de petites valeurs
avec de rares transactions à montant très élevé.

\begin{table}[H]
\centering
\caption{Statistiques descriptives de la variable \texttt{amount}}
\label{tab:amount_stats}
\begin{tabular}{ll}
\toprule
\textbf{Statistique} & \textbf{Valeur} \\
\midrule
Minimum     & 0,00 \\
Médiane     & $\sim$ 74~000 \\
Moyenne     & $\sim$ 179~862 \\
Maximum     & $\sim$ 92~445~517 \\
Skewness    & 30,80 \\
\bottomrule
\end{tabular}
\end{table}

Cette propriété justifie l'application d'une transformation logarithmique
$\log(1 + x)$ lors de la préparation des données pour corriger la skewness et
améliorer les performances des modèles linéaires.

\subsection{Analyse temporelle}

L'analyse de la variable \texttt{step} révèle une \textbf{saisonnalité des fraudes}.
Les heures à risque élevé identifiées dans l'EDA sont :
$\{0, 1, 2, 3, 4, 5, 6, 7, 8, 9, 23\}$ --- soit les heures nocturnes et
de début de matinée. Le ratio de fraude durant ces heures est de \textbf{10,63 fois
supérieur} à celui des heures ordinaires. Cette observation conduit à la création d'une
variable indicatrice \texttt{high\_risk\_hour}.

\subsection{Analyse des soldes et détection de leakage}

L'analyse des variables de soldes révèle des incohérences pour les transactions frauduleuses :
\texttt{oldbalanceOrg}, \texttt{newbalanceOrig}, \texttt{oldbalanceDest} et
\texttt{newbalanceDest} présentent des soldes nuls anormaux. Ces variables sont des
\textbf{indicateurs trop directs de la fraude} et pourraient causer une fuite
d'information (\textit{data leakage}) si elles étaient directement incluses dans les
modèles. À la place, une variable dérivée \texttt{balance\_diff\_orig} est créée, avec
une corrélation de \textbf{0,3662} avec la variable cible \texttt{isFraud}.

\subsection{Baseline du système existant}

Le système existant (\texttt{isFlaggedFraud}) sert de référence basse. Ses performances
illustrent l'inadéquation des approches par règles simples :

\begin{table}[H]
\centering
\caption{Performances de la baseline \texttt{isFlaggedFraud}}
\label{tab:baseline_perf}
\begin{tabular}{lr}
\toprule
\textbf{Métrique} & \textbf{Valeur} \\
\midrule
Recall    & 0,0039 \\
Precision & 1,0000 \\
F1-Score  & 0,0077 \\
\bottomrule
\end{tabular}
\end{table}

\textit{Ce système ne détecte qu'une infime fraction des fraudes réelles --- une seule
transaction sur 258 ---, confirmant la nécessité d'une approche par apprentissage
automatique.}

% ---------------------------------------------------
\section{Étude de l'existant et état de l'art}
% ---------------------------------------------------

\subsection{Approches supervisées}

Les approches supervisées nécessitent des étiquettes (\textit{labels}) pour chaque
transaction. Parmi les plus utilisées pour la détection de fraude :

\begin{itemize}
  \item \textbf{Régression logistique} : baseline linéaire, faible coût computationnel,
    bonne interprétabilité.
  \item \textbf{Forêt aléatoire} : ensembliste, robuste au bruit, gère bien les
    interactions non linéaires entre variables.
  \item \textbf{Gradient Boosting (XGBoost, LightGBM)} : performances état de l'art
    sur données tabulaires.
  \item \textbf{Réseaux de neurones feedforward} : puissants mais nécessitent beaucoup
    de données labellisées.
\end{itemize}

\subsection{Approches non supervisées}

Lorsque les étiquettes sont rares ou indisponibles, les approches non supervisées
permettent de détecter des patterns inhabituels sans connaissance préalable des fraudes :

\begin{itemize}
  \item \textbf{Isolation Forest} : isole les anomalies par partitions aléatoires
    successives.
  \item \textbf{One-Class SVM} : délimite la frontière de la normalité dans l'espace
    des features.
  \item \textbf{Autoencoder} : apprend à reconstruire les données normales ; une erreur
    de reconstruction élevée signale une anomalie.
  \item \textbf{Autoencoders variationnels (VAE)} : modélisation probabiliste de la
    distribution normale, plus robuste pour les données continues.
\end{itemize}

\subsection{Approches retenues pour ce projet}

Pour ce projet, trois approches complémentaires ont été retenues, couvrant le spectre
supervisé / non supervisé :

\begin{enumerate}
  \item \textbf{Régression logistique} --- baseline supervisée simple, pour établir
    une référence linéaire solide.
  \item \textbf{Forêt aléatoire} --- référence supervisée robuste, capable de capturer
    les interactions non linéaires entre variables.
  \item \textbf{Autoencoder} --- approche non supervisée orientée reconstruction, utile
    lorsque les labels sont limités ou lorsqu'on souhaite modéliser les patterns normaux
    sans dépendre des exemples de fraude.
\end{enumerate}

Ce choix vise à couvrir plusieurs angles d'attaque du problème, afin de disposer d'une
comparaison objective et d'une complémentarité méthodologique.

% ---------------------------------------------------
\section*{Conclusion}
\addcontentsline{toc}{section}{Conclusion}
% ---------------------------------------------------

L'analyse exploratoire a mis en évidence les caractéristiques clés du jeu de données :
déséquilibre extrême des classes (1:774), asymétrie des montants (skewness = 30,80),
saisonnalité temporelle des fraudes (ratio 10,63× la nuit) et risque de fuite d'information
sur les variables de soldes. Ces observations ne sont pas de simples curiosités statistiques
--- elles orientent directement chaque décision de préparation et de modélisation décrite
dans les chapitres suivants. L'inadéquation flagrante du système de règles existant
(F1 = 0,0077) confirme l'utilité d'une approche par machine learning. Le chapitre suivant
décrit en détail les opérations de préparation des données.

% =====================================================
% CHAPITRE 4 : PRÉPARATION DES DONNÉES (CRISP-DM Phase 3)
% =====================================================
\chapter{Préparation des données}

% ---------------------------------------------------
\section*{Introduction}
\addcontentsline{toc}{section}{Introduction}
% ---------------------------------------------------

Préparer les données, c'est poser les fondations sur lesquelles reposent tous les modèles
à venir. J'ai rapidement réalisé que les décisions prises ici auraient des répercussions
directes sur la validité de l'évaluation finale. L'enjeu central n'était pas seulement
technique, mais méthodologique : comment extraire des signaux pertinents des variables
brutes, tout en évitant d'introduire sans le vouloir de l'information du futur dans
l'apprentissage --- ce que l'on appelle le \textit{data leakage} ? Ce chapitre détaille
chaque choix effectué et les justifications empiriques qui les sous-tendent, telles
qu'elles ressortent des notebooks \texttt{01\_data\_understanding.ipynb} et
\texttt{02\_data\_preparation.ipynb}.

% ---------------------------------------------------
\section{Pipeline de préparation}
% ---------------------------------------------------

\subsection{Vue d'ensemble}

Le pipeline de préparation suit les étapes définies dans le module
\texttt{src/preprocessing/preprocessing.py}. Ces étapes sont appliquées dans un ordre
strict pour garantir l'étanchéité des partitions et l'absence de fuite d'information :

\begin{enumerate}
  \item Suppression des variables à risque de leakage
  \item Ingénierie de variables (\textit{feature engineering})
  \item Division stratifiée du corpus (train / validation / test)
  \item Normalisation (StandardScaler, \textit{fit} sur train uniquement)
  \item Gestion du déséquilibre (pondération et SMOTE)
\end{enumerate}

\subsection{Suppression des variables à risque de leakage}

La première décision consiste à identifier et éliminer les variables qui introduiraient de
l'information future ou qui révèleraient directement la classe cible. Ces suppressions sont
motivées empiriquement par l'analyse exploratoire :

\begin{table}[H]
\centering
\caption{Variables supprimées et justifications}
\label{tab:cols_dropped}
\begin{tabular}{p{4cm}p{8cm}}
\toprule
\textbf{Variable} & \textbf{Justification} \\
\midrule
\texttt{nameOrig}, \texttt{nameDest}
  & Identifiants nominaux non encodables sans risque de sur-apprentissage \\[0.2cm]
\texttt{oldbalanceOrg}, \texttt{newbalanceOrig}
  & Corrélation directe avec la fraude ; fuite d'information avérée \\[0.2cm]
\texttt{oldbalanceDest}, \texttt{newbalanceDest}
  & Idem ; révèlent les patterns frauduleux de manière trop évidente \\[0.2cm]
\texttt{isFlaggedFraud}
  & Label proxy du système existant ; ne doit pas alimenter l'apprentissage \\
\bottomrule
\end{tabular}
\end{table}

% ---------------------------------------------------
\section{Ingénierie de variables}
% ---------------------------------------------------

\subsection{Variables temporelles}

La variable \texttt{step} (pas de temps en heures) est décomposée en trois variables
cycliques qui capturent des patterns à différentes granularités temporelles :

\begin{itemize}
  \item \texttt{hour} : heure dans la journée (0--23), calculée par
    $\texttt{step} \bmod 24$ ;
  \item \texttt{day} : jour de simulation (0--30), calculé par
    $\lfloor \texttt{step} / 24 \rfloor \bmod 31$ ;
  \item \texttt{week} : semaine de simulation (0--4), calculée par
    $\lfloor \texttt{step} / 168 \rfloor$.
\end{itemize}

\subsection{Variables comportementales et de risque}

Cinq variables dérivées sont construites à partir des observations de l'EDA :

\begin{table}[H]
\centering
\caption{Variables dérivées construites lors du feature engineering}
\label{tab:features_derivees}
\begin{tabular}{p{3.5cm}p{4.5cm}p{4cm}}
\toprule
\textbf{Variable} & \textbf{Formule / Logique} & \textbf{Signal mesuré} \\
\midrule
\texttt{high\_risk\_hour} &
  $1$ si \texttt{hour} $\in \{0,...,9, 23\}$ &
  Ratio fraude $10{,}63\times$ supérieur \\[0.2cm]
\texttt{is\_transfer\_or\_cashout} &
  $1$ si \texttt{type} $\in$ \{TRANSFER, CASH\_OUT\} &
  Seuls types avec fraudes réelles \\[0.2cm]
\texttt{balance\_diff\_orig} &
  \texttt{oldbalanceOrg} $-$ \texttt{newbalanceOrig} &
  Corrélation 0{,}3662 avec \texttt{isFraud} \\[0.2cm]
\texttt{dest\_zero\_balance} &
  $1$ si \texttt{oldbalanceDest} $= 0$ &
  Taux fraude 2{,}52\,\% vs 0{,}13\,\% global \\[0.2cm]
\texttt{log\_amount} &
  $\log(1 + \texttt{amount})$ &
  Corrige skewness $= 30{,}80$ \\
\bottomrule
\end{tabular}
\end{table}

Chacune de ces variables est justifiée empiriquement par l'EDA --- aucun choix n'est
arbitraire.

\subsection{Encodage one-hot}

La variable catégorielle \texttt{type} (5 modalités) est encodée par
\textbf{one-hot encoding} (\textit{drop\_first=False}), produisant les colonnes
\texttt{type\_CASH\_IN}, \texttt{type\_CASH\_OUT}, \texttt{type\_DEBIT},
\texttt{type\_PAYMENT} et \texttt{type\_TRANSFER}. L'absence de suppression d'une
colonne (\textit{drop\_first=False}) est volontaire : avec un arbre de décision ou un
autoencoder, toutes les modalités sont nécessaires.

\subsection{Ensemble final des variables}

Après feature engineering, le jeu de données comprend \textbf{14 variables} :

\begin{table}[H]
\centering
\caption{Ensemble final des 14 variables du modèle}
\label{tab:final_features}
\begin{tabular}{p{2.5cm}p{6cm}p{3cm}}
\toprule
\textbf{Catégorie} & \textbf{Variables} & \textbf{Traitement} \\
\midrule
Temporelles  & \texttt{step}, \texttt{hour}, \texttt{day}, \texttt{week} & StandardScaler \\
Montant      & \texttt{log\_amount} & StandardScaler \\
Comportement & \texttt{balance\_diff\_orig} & StandardScaler \\
Binaires     & \texttt{high\_risk\_hour}, \texttt{is\_transfer\_or\_cashout},
               \texttt{dest\_zero\_balance} & --- \\
Type (OHE)   & \texttt{type\_CASH\_IN}, \texttt{type\_CASH\_OUT}, \texttt{type\_DEBIT},
               \texttt{type\_PAYMENT}, \texttt{type\_TRANSFER} & --- \\
\bottomrule
\end{tabular}
\end{table}

% ---------------------------------------------------
\section{Division et normalisation}
% ---------------------------------------------------

\subsection{Division stratifiée}

Le corpus est divisé en trois sous-ensembles de manière \textbf{stratifiée} sur la variable
cible, afin de conserver la proportion de fraudes dans chaque partition --- essentiel
lorsque les classes sont aussi déséquilibrées.

\begin{table}[H]
\centering
\caption{Division stratifiée du corpus}
\label{tab:splits}
\begin{tabular}{lrrl}
\toprule
\textbf{Ensemble} & \textbf{Taille} & \textbf{Fraudes} & \textbf{Rôle} \\
\midrule
Entraînement & 139~999 & 181 & Apprentissage des modèles \\
Validation   & 30~000  & 38  & Réglage des hyperparamètres et seuils \\
Test         & 30~001  & 39  & Évaluation finale (jamais consulté avant) \\
\bottomrule
\end{tabular}
\end{table}

\textbf{Principe d'étanchéité :} L'ensemble de test n'est jamais consulté avant
l'évaluation finale, conformément aux bonnes pratiques anti-\textit{data snooping}.
Ce principe a été respecté tout au long du projet.

\subsection{Normalisation}

Le \textbf{StandardScaler} est ajusté (\textit{fit}) exclusivement sur l'ensemble
d'entraînement, puis appliqué (\textit{transform}) sur les trois partitions. Cette
discipline stricte garantit l'absence de leakage statistique : aucune information sur la
distribution du test ne s'infiltre dans les transformations appliquées à l'entraînement.

% ---------------------------------------------------
\section{Gestion du déséquilibre}
% ---------------------------------------------------

\subsection{Pondération des classes}

Pour les modèles supervisés, des poids de classes sont calculés automatiquement
(formule \textit{sklearn balanced}) :

\[
w_c = \frac{n_{\text{total}}}{n_{\text{classes}} \times n_c}
\]

Avec 139~999 transactions d'entraînement dont 181 fraudes, les poids obtenus sont
approximativement $w_0 \approx 0{,}50$ et $w_1 \approx 386{,}7$. Cela revient à donner
à chaque transaction frauduleuse un poids environ 386 fois supérieur à une transaction
légitime lors du calcul de la fonction de perte --- compensant ainsi le déséquilibre du
corpus.

\subsection{SMOTE}

Pour les variantes avec rééchantillonnage, \textbf{SMOTE} (\textit{Synthetic Minority
Over-sampling Technique}) génère des exemples synthétiques de la classe minoritaire par
interpolation dans l'espace des features. Le principe : pour chaque transaction frauduleuse,
SMOTE choisit aléatoirement $k$ voisins les plus proches dans l'espace des features, puis
crée de nouveaux exemples sur les segments reliant la transaction à ses voisins.

Il est appliqué \textbf{uniquement sur l'ensemble d'entraînement} (jamais sur validation
ou test), afin de ne pas contaminer l'évaluation avec des exemples synthétiques.

\subsection{Ensemble pour l'autoencoder}

Pour l'autoencoder non supervisé, un sous-ensemble \texttt{X\_train\_normal} est extrait :
il contient \textbf{uniquement les 139~818 transactions légitimes} de l'entraînement
(0 fraude). Cela permet au modèle d'apprendre exclusivement la distribution normale ---
principe fondateur de la détection d'anomalies par reconstruction.

% ---------------------------------------------------
\section*{Conclusion}
\addcontentsline{toc}{section}{Conclusion}
% ---------------------------------------------------

La préparation des données produit un jeu de \textbf{14 variables représentatives} des
patterns de fraude, en respectant strictement l'étanchéité entre les partitions. Les
7 variables dérivées ont chacune été validées par une mesure empirique issue de l'EDA ---
aucune n'est arbitraire. La discipline anti-leakage (StandardScaler fit sur train seulement,
SMOTE appliqué uniquement sur le train, test jamais consulté) garantit que l'évaluation
finale reflète les véritables performances sur des données inédites. Ce socle solide permet
d'entraîner des modèles fiables et comparables, décrits dans le chapitre suivant.

% =====================================================
% CHAPITRE 5 : MODÉLISATION (CRISP-DM Phase 4)
% =====================================================
\chapter{Modélisation}

% ---------------------------------------------------
\section*{Introduction}
\addcontentsline{toc}{section}{Introduction}
% ---------------------------------------------------

Une fois les données préparées, vient la phase de construction des modèles. Plutôt que de
choisir d'emblée l'approche la plus sophistiquée, j'ai opté pour une progression logique :
commencer par un modèle simple servant de référence, complexifier si les performances le
justifient, puis compléter avec une perspective non supervisée. Cette stratégie permet
d'évaluer la valeur ajoutée réelle de chaque niveau de complexité et d'éviter un
sur-ingénierie inutile. Ce chapitre décrit les trois architectures développées et les
décisions qui ont guidé leur conception, à travers les notebooks
\texttt{03\_baseline\_models.ipynb} et \texttt{04\_autoencoder.ipynb}.

% ---------------------------------------------------
\section{Régression logistique}
% ---------------------------------------------------

\subsection{Principe et rôle dans le projet}

La régression logistique est la \textbf{baseline} --- le modèle de référence que tout
algorithme plus complexe doit surpasser pour justifier sa complexité. Elle produit une
probabilité d'appartenance à la classe fraude via la fonction sigmoïde :

\[
P(\text{fraude} \mid \mathbf{x}) = \sigma(\mathbf{w}^\top \mathbf{x} + b)
= \frac{1}{1 + e^{-(\mathbf{w}^\top \mathbf{x} + b)}}
\]

Sa faible capacité expressive est ici un avantage : si la régression logistique détecte
déjà une grande partie des fraudes, cela signifie que le signal est linéairement séparable,
ce qui simplifie l'ensemble de la chaîne. Dans le cas contraire, cela justifie le passage
à des modèles non linéaires.

\subsection{Hyperparamètres et configurations évaluées}

\begin{table}[H]
\centering
\caption{Hyperparamètres de la régression logistique}
\label{tab:lr_params}
\begin{tabular}{lll}
\toprule
\textbf{Paramètre} & \textbf{Valeur} & \textbf{Justification} \\
\midrule
\texttt{C}             & 0,1      & Régularisation forte (peu de fraudes) \\
\texttt{solver}        & lbfgs   & Adapté aux datasets de taille moyenne \\
\texttt{class\_weight} & balanced & Correction automatique du déséquilibre \\
\texttt{max\_iter}     & 1000    & Assure la convergence \\
\bottomrule
\end{tabular}
\end{table}

Deux configurations sont comparées :
\begin{itemize}
  \item \textbf{LR\_balanced} : entraînée sur \texttt{X\_train} avec
    \texttt{class\_weight='balanced'} ;
  \item \textbf{LR\_smote} : entraînée sur \texttt{X\_smote} (données équilibrées par
    SMOTE) sans pondération explicite.
\end{itemize}

\subsection{Réglage du seuil de décision}

À seuil par défaut ($0{,}5$), \texttt{LR\_balanced} maximise fortement le rappel mais au
prix d'une précision très faible. Une recherche de seuil sur l'ensemble de validation est
donc réalisée à partir des probabilités produites par le modèle. Les seuils optimaux retenus
sont \textbf{0,9986} pour \texttt{LR\_balanced} et \textbf{0,9621} pour \texttt{LR\_smote}.
Après ce recalibrage, \texttt{LR\_smote} devient la meilleure variante logistique avec un
rappel de \textbf{0,6923}, une précision de \textbf{0,7297} et un F1-score de
\textbf{0,7105} sur l'ensemble de test.

% ---------------------------------------------------
\section{Forêt aléatoire}
% ---------------------------------------------------

\subsection{Principe et avantages}

La forêt aléatoire constitue la \textbf{référence supervisée robuste} de ce projet. Son
principe repose sur la diversité : entraîner un grand nombre d'arbres de décision sur des
sous-échantillons distincts des données (\textit{bootstrap}), puis agréger leurs
prédictions. Cette diversité rend l'ensemble résistant au sur-apprentissage
(\textit{overfitting}) et capable de capturer des interactions non linéaires entre variables
que la régression logistique ne peut pas modéliser.

La prédiction finale est la \textbf{moyenne des probabilités} émises par chaque arbre
individuel, ce qui produit des scores bien calibrés.

\subsection{Hyperparamètres}

\begin{table}[H]
\centering
\caption{Hyperparamètres de la forêt aléatoire}
\label{tab:rf_params}
\begin{tabular}{lll}
\toprule
\textbf{Paramètre} & \textbf{Valeur} & \textbf{Justification} \\
\midrule
\texttt{n\_estimators}      & 300 & Stabilité des estimations \\
\texttt{max\_depth}         & 10  & Limite l'overfitting sur peu de fraudes \\
\texttt{min\_samples\_leaf} & 5   & Évite les feuilles trop petites \\
\texttt{class\_weight}      & balanced & Correction du déséquilibre \\
\bottomrule
\end{tabular}
\end{table}

Deux variantes sont comparées :
\begin{itemize}
  \item \textbf{RF\_balanced} : entraînée sur \texttt{X\_train} avec
    \texttt{class\_weight='balanced'} ;
  \item \textbf{RF\_smote} : entraînée sur \texttt{X\_smote}, sans pondération explicite.
\end{itemize}

\subsection{Réglage du seuil et importance des variables}

Les seuils optimaux retenus après recherche sur validation sont \textbf{0,6235} pour
\texttt{RF\_balanced} et \textbf{0,6291} pour \texttt{RF\_smote}. La variante
\texttt{RF\_smote} obtient les meilleures performances supervisées globales sur le test :
rappel de \textbf{0,7949}, précision de \textbf{0,8158} et F1-score de \textbf{0,8052}.

La forêt aléatoire permet également de calculer l'importance de chaque variable via la
réduction moyenne d'impureté (\textit{Mean Decrease in Impurity}). Sur
\texttt{RF\_smote}, les cinq variables les plus influentes sont :

\begin{table}[H]
\centering
\caption{Importance des variables (RF\_smote)}
\label{tab:feature_importance_rf}
\begin{tabular}{lr}
\toprule
\textbf{Variable} & \textbf{Importance} \\
\midrule
\texttt{balance\_diff\_orig}       & $\sim$0,385 \\
\texttt{hour}                      & $\sim$0,118 \\
\texttt{dest\_zero\_balance}       & $\sim$0,115 \\
\texttt{type\_TRANSFER}            & $\sim$0,089 \\
\texttt{log\_amount}               & $\sim$0,081 \\
\bottomrule
\end{tabular}
\end{table}

Cette hiérarchie confirme que le comportement du compte émetteur, la temporalité et la
nature du virement sont les signaux les plus discriminants --- en cohérence avec l'EDA.

% ---------------------------------------------------
\section{Autoencoder}
% ---------------------------------------------------

\subsection{Principe de détection par reconstruction}

L'autoencoder repose sur une idée intuitive : \textbf{apprendre à compresser puis à
reconstruire fidèlement les données normales}. Si une transaction légitime peut être
reproduite depuis un espace latent de petite dimension, c'est qu'elle s'inscrit dans des
patterns réguliers. Une transaction frauduleuse --- jamais vue à l'entraînement --- sera
en revanche mal reconstruite, trahissant son caractère inhabituel par une erreur de
reconstruction élevée.

\begin{definitionbox}
\textbf{Principe de l'autoencoder pour la détection d'anomalies :}
\begin{itemize}
  \item \textbf{Entraînement :} uniquement sur les transactions légitimes → le modèle
    apprend à reconstruire fidèlement les patterns normaux.
  \item \textbf{Inférence :} sur une transaction frauduleuse (pattern jamais vu),
    l'erreur de reconstruction est élevée car le modèle ne sait pas comment la reproduire.
  \item \textbf{Décision :} si l'erreur de reconstruction (MSE) dépasse un seuil
    $\theta$, la transaction est classée comme anomalie.
\end{itemize}
\end{definitionbox}

\subsection{Architecture du réseau}

L'architecture retenue est un autoencoder symétrique à trois couches dans l'encodeur et
trois dans le décodeur, avec un espace latent (\textit{bottleneck}) de dimension 4 :

\begin{table}[H]
\centering
\caption{Architecture de l'autoencoder ($14 \to 10 \to 7 \to [4] \to 7 \to 10 \to 14$)}
\label{tab:ae_architecture}
\begin{tabular}{llll}
\toprule
\textbf{Partie} & \textbf{Couche} & \textbf{Dimension} & \textbf{Activation} \\
\midrule
Entrée   & Input      & 14 & --- \\
Encodeur & Dense 1    & 10 & ReLU + BatchNorm + Dropout(0,2) \\
         & Dense 2    & 7  & ReLU + BatchNorm + Dropout(0,2) \\
         & Bottleneck & 4  & ReLU \\
Décodeur & Dense 3    & 7  & ReLU + BatchNorm + Dropout(0,2) \\
         & Dense 4    & 10 & ReLU + BatchNorm + Dropout(0,2) \\
Sortie   & Dense 5    & 14 & Linéaire \\
\bottomrule
\end{tabular}
\end{table}

Les techniques de régularisation employées sont :
\begin{itemize}
  \item \textbf{BatchNorm} (normalisation par lots) : stabilise l'apprentissage ;
  \item \textbf{Dropout} (taux 0,2) : réduit le sur-apprentissage ;
  \item \textbf{Régularisation L2} ($\lambda = 10^{-5}$) sur les poids denses.
\end{itemize}

\subsection{Entraînement}

\begin{table}[H]
\centering
\caption{Paramètres d'entraînement de l'autoencoder}
\label{tab:ae_training}
\begin{tabular}{ll}
\toprule
\textbf{Paramètre} & \textbf{Valeur} \\
\midrule
Données d'entraînement  & 139~818 transactions légitimes (0 fraude) \\
Validation interne      & 10\,\% du train normal \\
Fonction de perte       & MSE (Mean Squared Error) \\
Optimiseur              & Adam ($\alpha = 10^{-3}$) \\
Batch size              & 256 \\
Epochs max              & 100 (EarlyStopping, patience = 10) \\
ReduceLROnPlateau       & facteur 0,5, patience 5, min\_lr $10^{-6}$ \\
\bottomrule
\end{tabular}
\end{table}

L'entraînement s'arrête après \textbf{35 époques effectives} grâce à \textit{EarlyStopping}.
Le meilleur \texttt{val\_loss} observé est de \textbf{0,145024} et le temps d'entraînement
est de \textbf{45,2 secondes}. Le modèle comporte \textbf{664 paramètres} entraînables,
ce qui reste très léger au regard de la taille du corpus.

\subsection{Détermination du seuil de détection}

Une première borne est estimée par le \textbf{95\textsuperscript{e} percentile} de
l'erreur de reconstruction sur les transactions normales d'entraînement
($\theta_{p95} = 0{,}403750$). Ce seuil est utile comme repère initial, mais trop permissif
pour une décision finale.

Le seuil optimal est ensuite trouvé sur l'ensemble de validation (\texttt{X\_val}) en
balayant 200 valeurs entre le 50\textsuperscript{e} et le 99,9\textsuperscript{e}
percentile des erreurs de reconstruction. Pour chaque valeur, le F1-score est calculé.
Le seuil maximisant ce score est retenu : \textbf{$\theta = 1{,}687408$}.

\subsection{Visualisation de l'espace latent}

La projection PCA du bottleneck (dimension 4) en 2D permet de vérifier qualitativement
que les fraudes de validation tendent à se positionner \textbf{en périphérie des zones de
forte densité normale} --- confirmant la pertinence du mécanisme de reconstruction comme
score d'anomalie.

% ---------------------------------------------------
\section{Récapitulatif des configurations}
% ---------------------------------------------------

\begin{table}[H]
\centering
\caption{Récapitulatif des configurations de modèles construites}
\label{tab:models_recap}
\begin{tabular}{p{2.8cm}p{3cm}p{3cm}p{3.5cm}}
\toprule
\textbf{Configuration} & \textbf{Type} & \textbf{Données train}
  & \textbf{Gestion déséquilibre} \\
\midrule
LR\_balanced  & Supervisé linéaire   & \texttt{X\_train}        & \texttt{class\_weight} \\
LR\_smote     & Supervisé linéaire   & \texttt{X\_smote}        & SMOTE \\
RF\_balanced  & Supervisé ensembliste & \texttt{X\_train}       & \texttt{class\_weight} \\
RF\_smote     & Supervisé ensembliste & \texttt{X\_smote}       & SMOTE \\
Autoencoder   & Non supervisé        & \texttt{X\_train\_normal} & Données normales seules \\
\bottomrule
\end{tabular}
\end{table}

% ---------------------------------------------------
\section*{Conclusion}
\addcontentsline{toc}{section}{Conclusion}
% ---------------------------------------------------

Ce chapitre a décrit la conception et l'implémentation des cinq configurations de modèles
étudiées. Les variantes supervisées montrent que le choix du seuil de décision influence
fortement le compromis précision/rappel, et que \texttt{RF\_smote} s'impose comme la
meilleure configuration supervisée. L'autoencoder complète ce dispositif par une logique
non supervisée fondée sur la reconstruction, particulièrement utile lorsque les labels
sont rares. Les premiers résultats sur l'ensemble de validation suggèrent que la forêt
aléatoire dépasse les objectifs fixés (recall $> 0{,}70$, F1 $> 0{,}70$). L'évaluation
comparative détaillée sur l'ensemble de test, ainsi que les mécanismes d'explicabilité,
sont présentés dans le chapitre suivant.

% =====================================================
% CHAPITRE 6 : ÉVALUATION ET INTERPRÉTABILITÉ (CRISP-DM Phase 5)
% =====================================================
\chapter{Évaluation et interprétabilité}

% ---------------------------------------------------
\section*{Introduction}
\addcontentsline{toc}{section}{Introduction}
% ---------------------------------------------------

Avoir un modèle qui fonctionne sur les données d'entraînement ne suffit pas. L'évaluation
sert à répondre à deux questions distinctes : le modèle détecte-t-il réellement les fraudes
(évaluation quantitative sur métriques adaptées au déséquilibre) ? Et ses décisions
sont-elles compréhensibles et défendables pour un auditeur (évaluation qualitative via
SHAP, LIME et Ollama) ? Ces deux dimensions sont indissociables dans un contexte professionnel
où les alertes doivent non seulement être justes, mais aussi justifiées. Ce chapitre
synthétise les résultats issus des notebooks \texttt{03\_baseline\_models.ipynb},
\texttt{04\_autoencoder.ipynb}, \texttt{05\_ollama\_integration.ipynb} et
\texttt{06\_shap\_lime.ipynb}.

% ---------------------------------------------------
\section{Métriques d'évaluation}
% ---------------------------------------------------

\subsection{Choix et justification des métriques}

Dans un contexte de fort déséquilibre, l'exactitude (\textit{accuracy}) n'est pas
pertinente : un modèle qui prédit systématiquement « légitime » obtiendrait 99,87\,\%
d'exactitude tout en ne détectant aucune fraude. Les métriques retenues sont :

\begin{itemize}
  \item \textbf{Rappel (Recall)} : fraction des fraudes réelles correctement détectées.
    \textit{Priorité~1} car manquer une fraude a un coût financier et réputationnel direct.
    \[
    \text{Recall} = \frac{TP}{TP + FN}
    \]
  \item \textbf{Précision (Precision)} : fraction des alertes qui sont réellement des
    fraudes. \textit{Priorité~2} car trop de faux positifs surchargent les équipes
    d'investigation.
    \[
    \text{Precision} = \frac{TP}{TP + FP}
    \]
  \item \textbf{F1-Score} : moyenne harmonique de la précision et du rappel. Synthèse
    unique du compromis.
    \[
    \text{F1} = \frac{2 \cdot \text{Precision} \cdot \text{Recall}}
    {\text{Precision} + \text{Recall}}
    \]
  \item \textbf{PR-AUC} (\textit{Area Under the Precision-Recall Curve}) : synthétise
    le tradeoff précision/rappel sur l'ensemble des seuils possibles --- robuste au
    déséquilibre des classes.
\end{itemize}

\subsection{Résultats comparatifs sur l'ensemble de test}

\begin{table}[H]
\centering
\caption{Comparaison des performances sur l'ensemble de test}
\label{tab:results_comparison}
\begin{tabular}{lccccc}
\toprule
\textbf{Modèle} & \textbf{Seuil} & \textbf{Recall}
  & \textbf{Precision} & \textbf{F1-Score} & \textbf{PR-AUC} \\
\midrule
Baseline (\texttt{isFlaggedFraud}) & N/A    & 0,0039 & 1,0000 & 0,0077 & N/A \\
LR\_balanced                       & 0,9986 & 0,6410 & 0,6579 & 0,6494 & 0,7044 \\
LR\_smote                          & 0,9621 & 0,6923 & 0,7297 & 0,7105 & 0,7393 \\
RF\_balanced                       & 0,6235 & 0,7692 & 0,8333 & 0,8000 & 0,7794 \\
RF\_smote                          & 0,6291 & 0,7949 & 0,8158 & 0,8052 & 0,8405 \\
Autoencoder                        & 1,6874 & 0,3590 & 0,6087 & 0,4516 & 0,3734 \\
\bottomrule
\end{tabular}
\end{table}

\textit{Note : les seuils sont sélectionnés sur la validation, puis figés pour l'évaluation
finale sur l'ensemble de test --- jamais consulté lors de la phase d'optimisation.}

\subsection{Analyse de la matrice de confusion}

La matrice de confusion permet d'analyser précisément les types d'erreurs :
\begin{itemize}
  \item \textbf{Faux négatifs (FN)} : fraudes manquées --- coût financier direct,
    impact réputationnel, responsabilité de l'auditeur.
  \item \textbf{Faux positifs (FP)} : alertes injustifiées --- surcharge des équipes
    d'investigation, coût opérationnel.
\end{itemize}

La forêt aléatoire avec SMOTE offre le meilleur équilibre sur l'ensemble de test :
\textbf{31 fraudes détectées sur 39} pour seulement \textbf{7 faux positifs}. La variante
\texttt{RF\_balanced} reste très proche (30 TP, 6 FP), tandis que \texttt{LR\_smote}
conserve un compromis acceptable (27 TP, 10 FP). L'autoencoder, en revanche, ne détecte
que \textbf{14 fraudes sur 39} avec \textbf{9 faux positifs} --- moins performant que les
approches supervisées, mais conservant un intérêt lorsque les labels sont rares ou
indisponibles.

\subsection{Discussion et sélection du modèle}

Tous les critères de succès définis en phase~1 sont atteints pour la forêt aléatoire avec
SMOTE :

\begin{table}[H]
\centering
\caption{Vérification des critères de succès pour RF\_smote}
\label{tab:criteres_verification}
\begin{tabular}{llll}
\toprule
\textbf{Critère} & \textbf{Cible} & \textbf{Résultat} & \textbf{Statut} \\
\midrule
Recall   & $> 0{,}70$ & 0,7949 & \checkmark Atteint \\
F1-Score & $> 0{,}70$ & 0,8052 & \checkmark Atteint \\
PR-AUC   & $> 0{,}75$ & 0,8405 & \checkmark Atteint \\
\bottomrule
\end{tabular}
\end{table}

% ---------------------------------------------------
\section{Interprétabilité des modèles}
% ---------------------------------------------------

\subsection{Méthodes SHAP}

L'un des reproches fréquents adressés aux modèles d'apprentissage automatique est leur
opacité. Dans un contexte d'audit, un analyste ne peut se contenter d'une alerte : il lui
faut une justification claire des variables qui ont conduit le modèle à signaler une
transaction.

\begin{definitionbox}
\textbf{SHAP} (\textit{SHapley Additive exPlanations}) est une méthode d'interprétabilité
fondée sur la théorie des jeux coopératifs. Elle attribue à chaque variable une valeur
représentant sa contribution marginale à la prédiction, calculée en moyennant l'impact de
la variable sur toutes les combinaisons possibles de features. Cette approche garantit
additivité, cohérence et absence de biais entre les variables.
\end{definitionbox}

Pour la forêt aléatoire, un explainer \textbf{TreeSHAP} est utilisé --- plus rapide que
l'approche agnostique et exact pour les modèles arborescents. Les résultats montrent que
\texttt{balance\_diff\_orig}, \texttt{is\_transfer\_or\_cashout} et
\texttt{log\_amount} sont les variables à plus forte contribution, en cohérence avec
l'analyse exploratoire et l'importance des variables de la forêt.

\textbf{Deux niveaux d'analyse SHAP sont produits :}
\begin{itemize}
  \item \textbf{Global} : beeswarm plot montrant la distribution des valeurs SHAP sur
    l'ensemble du test, permettant d'identifier les variables globalement les plus
    influentes.
  \item \textbf{Local} : waterfall plot pour une transaction spécifique, montrant
    comment chaque feature pousse la prédiction vers « fraude » ou « légitime ».
\end{itemize}

\subsection{Méthodes LIME}

\begin{definitionbox}
\textbf{LIME} (\textit{Local Interpretable Model-agnostic Explanations}) explique chaque
prédiction individuellement en perturbant légèrement la transaction d'entrée et en
observant comment la prédiction varie. Il approxime localement le comportement du modèle
complexe par un modèle linéaire simple, interprétable par construction, valable uniquement
dans le voisinage de l'exemple analysé.
\end{definitionbox}

LIME est particulièrement utile pour expliquer un cas suspect spécifique à un analyste,
sans avoir besoin de comprendre le comportement global du modèle. Contrairement à SHAP,
il est \textbf{agnostique} à l'architecture du modèle et peut s'appliquer à n'importe
quel algorithme --- y compris l'autoencoder, pour lequel SHAP n'est pas nativement adapté.

\subsection{Comparaison SHAP vs LIME}

\begin{table}[H]
\centering
\caption{Comparaison des méthodes d'explicabilité}
\label{tab:shap_lime}
\begin{tabular}{lll}
\toprule
\textbf{Critère} & \textbf{SHAP} & \textbf{LIME} \\
\midrule
Scope          & Global + Local & Local uniquement \\
Fondement      & Théorie des jeux & Approximation locale linéaire \\
Fidélité       & Élevée (TreeSHAP exact) & Variable (approximation) \\
Rapidité       & Rapide (TreeSHAP) & Lent (nombreuses perturbations) \\
Agnosticisme   & Méthodes spécifiques disponibles & Entièrement agnostique \\
\bottomrule
\end{tabular}
\end{table}

% ---------------------------------------------------
\section{Intégration Ollama}
% ---------------------------------------------------

\subsection{Pourquoi un LLM local ?}

Ollama est un outil permettant d'exécuter des grands modèles de langage (LLM)
\textbf{entièrement en local}, sans envoyer de données vers des services cloud externes.
Ce point était non négociable dans le cadre de ce projet : les transactions financières
auditées sont confidentielles et ne peuvent pas quitter l'environnement sécurisé du client.
L'intégration avec Ollama permet de produire des explications en langage naturel ---
directement lisibles par un auditeur sans formation en machine learning --- sans compromettre
la confidentialité des données.

\subsection{Architecture de l'intégration}

Le module \texttt{src/llm\_integration/llm\_helper.py} encapsule la logique d'appel au
modèle de langage. Pour chaque transaction signalée comme anormale, les éléments suivants
sont transmis au LLM dans un prompt structuré :

\begin{itemize}
  \item Le score d'anomalie ou d'erreur de reconstruction de l'autoencoder ;
  \item Les valeurs des variables les plus contributives (issues de SHAP/LIME) ;
  \item Le contexte métier de la transaction (type, montant, heure, différence de solde).
\end{itemize}

Le LLM génère une explication structurée en quatre points : type de transaction à risque,
heure suspecte, comportement anormal des soldes, et recommandation d'action.

\subsection{Difficultés rencontrées}

\begin{itemize}
  \item \textbf{Instabilité des seuils de l'autoencoder} : de légères variations du seuil
    peuvent significativement modifier le ratio précision/rappel, nécessitant une
    optimisation soigneuse sur la validation.
  \item \textbf{Variabilité des explications LLM} : les explications Ollama peuvent varier
    selon la version du modèle utilisé et la qualité du prompt ; un prompt engineering
    rigoureux a été nécessaire.
  \item \textbf{Temps de calcul SHAP} : le calcul des valeurs SHAP pour l'ensemble de
    test complet nécessite plusieurs minutes sur CPU, ce qui a motivé la sauvegarde des
    tenseurs SHAP dans \texttt{outputs/models/shap\_values\_*.npy}.
\end{itemize}

% ---------------------------------------------------
\section*{Conclusion}
\addcontentsline{toc}{section}{Conclusion}
% ---------------------------------------------------

L'évaluation confirme que la forêt aléatoire avec SMOTE constitue la \textbf{meilleure
approche supervisée} pour ce jeu de données, avec un recall de 0,7949, un F1-score de
0,8052 et un PR-AUC de 0,8405 --- tous au-dessus des cibles fixées. L'autoencoder apporte
une perspective complémentaire non supervisée, utile lorsque les étiquettes sont limitées.
La \textbf{double couche d'explicabilité} (SHAP/LIME + Ollama) répond à l'objectif
d'interprétabilité en rendant les prédictions compréhensibles tant pour les data scientists
que pour les analystes métier non techniciens. Le chapitre suivant présente les aspects
de déploiement et les garanties de reproductibilité du pipeline complet.

% =====================================================
% CHAPITRE 7 : DÉPLOIEMENT ET REPRODUCTIBILITÉ (CRISP-DM Phase 6)
% =====================================================
\chapter{Déploiement et reproductibilité}

% ---------------------------------------------------
\section*{Introduction}
\addcontentsline{toc}{section}{Introduction}
% ---------------------------------------------------

Un modèle performant qui ne peut pas être réexécuté par quelqu'un d'autre n'a pas grande
valeur dans un contexte professionnel. Pour ce projet, le déploiement signifie rendre
l'ensemble du travail \textbf{reproductible de bout en bout}, maintenable par un futur
contributeur et extensible si les besoins évoluent. Ce chapitre décrit l'architecture
choisie, l'organisation des sorties, le script d'orchestration, les garanties de
reproductibilité, la suite de tests unitaires mise en place, et les perspectives
d'amélioration identifiées pour une éventuelle industrialisation du système.

% ---------------------------------------------------
\section{Architecture du projet}
% ---------------------------------------------------

\subsection{Structure modulaire}

Le projet est organisé en modules fonctionnels clairement délimités, chacun ayant une
responsabilité unique (\textit{separation of concerns}) :

\begin{table}[H]
\centering
\caption{Structure des modules du projet}
\label{tab:modules}
\begin{tabular}{p{4.5cm}p{8cm}}
\toprule
\textbf{Module} & \textbf{Responsabilité} \\
\midrule
\texttt{src/preprocessing/}        & Chargement, nettoyage, normalisation, split \\
\texttt{src/feature\_engineering/} & Création des 7 variables dérivées \\
\texttt{src/models/}               & Modèles ML (LR, RF) et autoencoder PyTorch \\
\texttt{src/pipeline/}             & Orchestration entraînement et inférence \\
\texttt{src/utils/}                & Métriques, évaluation, utilitaires communs \\
\texttt{src/visualization/}        & Génération de figures et graphiques \\
\texttt{src/ollama\_integration/}  & Génération d'explications LLM via Ollama \\
\texttt{src/explainability/}       & Explicabilité SHAP et LIME \\
\texttt{tests/}                    & Tests unitaires (pytest) \\
\bottomrule
\end{tabular}
\end{table}

Cette organisation garantit que chaque module peut être modifié, testé et remplacé
indépendamment --- un avantage décisif pour la maintenance et l'évolution du système.

\subsection{Organisation des sorties}

Le pipeline produit des artefacts organisés dans le répertoire \texttt{outputs/} :

\begin{itemize}
  \item \texttt{outputs/models/} : artefacts sérialisés --- poids de l'autoencoder,
    scalers, seuils optimaux, poids des classes, scores AE
    (\texttt{ae\_scores\_val/test.npy}) et tenseurs SHAP sauvegardés
    (\texttt{shap\_values\_*.npy}) ;
  \item \texttt{outputs/reports/} : rapports JSON des phases exécutées
    (\texttt{eda\_report.json}, \texttt{prep\_report.json},
    \texttt{baseline\_report.json}, \texttt{autoencoder\_report.json}) ;
  \item \texttt{outputs/figures/} : courbes de performance, matrices de confusion,
    distributions, beeswarm SHAP, projection latente de l'AE ;
  \item \texttt{outputs/anomalies\_report.csv} : liste des transactions signalées avec
    leurs scores et niveaux de risque ;
  \item \texttt{outputs/explanations.json} : explications textuelles des anomalies
    générées par Ollama.
\end{itemize}

% ---------------------------------------------------
\section{Orchestration et reproductibilité}
% ---------------------------------------------------

\subsection{Script d'orchestration}

Le script \texttt{run\_all.py} constitue le point d'entrée unique du pipeline. Il valide
d'abord les imports des modules \texttt{src/}, puis exécute séquentiellement les six
notebooks (\texttt{01} à \texttt{06}) via \texttt{nbconvert}. Il propose plusieurs
modes d'exécution :

\begin{itemize}
  \item \textbf{Exécution complète} : lance les notebooks 01 à 06 en séquence ;
  \item \textbf{Reprise partielle} : option \texttt{--from N} pour reprendre à partir
    d'un notebook spécifique ;
  \item \textbf{Notebook unique} : option \texttt{--only N} pour n'exécuter qu'une seule
    étape ;
  \item \textbf{Vérification uniquement} : option \texttt{--check-only} pour valider
    les imports sans lancer les calculs.
\end{itemize}

\subsection{Garanties de reproductibilité}

La reproductibilité est assurée par quatre mécanismes complémentaires :

\begin{enumerate}
  \item \textbf{Graines aléatoires fixes} : \texttt{RANDOM\_STATE = 42} dans tous les
    modules qui comportent une part d'aléatoire (split, SMOTE, forêt aléatoire, dropout).
  \item \textbf{Environnement virtuel} : toutes les dépendances sont centralisées dans
    \texttt{requirements.txt} avec des contraintes de versions, à installer dans un
    environnement Python isolé.
  \item \textbf{Pipeline déterministe} : toutes les transformations (scaler, encodeurs)
    sont sérialisées via \texttt{joblib} dans \texttt{outputs/models/} pour être
    réappliquées à l'inférence sans recalcul.
  \item \textbf{Séparation stricte des données} : les ensembles de validation et test
    ne sont jamais utilisés pour l'apprentissage ni les transformations statistiques.
\end{enumerate}

% ---------------------------------------------------
\section{Tests unitaires}
% ---------------------------------------------------

\subsection{Structure de la suite de tests}

Le dépôt contient une suite de tests unitaires structurée sous \texttt{tests/} et
exécutable via \texttt{pytest -q} :

\begin{table}[H]
\centering
\caption{Suite de tests unitaires du projet}
\label{tab:tests}
\begin{tabular}{p{5.5cm}p{7cm}}
\toprule
\textbf{Fichier de test} & \textbf{Modules testés} \\
\midrule
\texttt{tests/test\_preprocessing.py}  & Pipeline de préparation des données \\
\texttt{tests/test\_models.py}         & Modèles ML et autoencoder \\
\texttt{tests/test\_visualization.py}  & Génération de figures \\
\texttt{tests/test\_utils.py}          & Utilitaires et métriques \\
\texttt{tests/test\_ollama.py}         & Intégration LLM (Ollama) \\
\bottomrule
\end{tabular}
\end{table}

Les fichiers de tests constituent un squelette structuré prêt à être alimenté. L'étape
suivante consiste à implémenter les cas de test dans ces fichiers pour atteindre
l'objectif de couverture à 80\,\%.

% ---------------------------------------------------
\section{Perspectives d'amélioration}
% ---------------------------------------------------

\subsection{Améliorations techniques}

Les résultats obtenus ouvrent plusieurs pistes d'amélioration algorithmique :

\begin{itemize}
  \item \textbf{Modèles de gradient boosting} : intégration de LightGBM ou XGBoost,
    qui constituent l'état de l'art sur les données tabulaires et pourraient améliorer
    les performances supervisées tout en restant interprétables via SHAP.
  \item \textbf{Autoencoder variationnel (VAE)} : passage d'un autoencoder déterministe
    à un VAE qui modélise la distribution des données normales sous forme probabiliste
    --- plus robuste pour détecter les anomalies de bord de distribution.
  \item \textbf{Détection en temps réel} : portage du pipeline vers une architecture
    streaming (Apache Kafka + Apache Flink) pour analyser les transactions au fil de
    l'eau, plutôt qu'en batch.
  \item \textbf{Apprentissage en ligne} : mise à jour continue des modèles au fil des
    nouvelles transactions validées par les auditeurs, pour s'adapter aux évolutions
    des comportements frauduleux.
\end{itemize}

\subsection{Améliorations métier}

Sur le plan fonctionnel, plusieurs extensions rendraient le système directement opérationnel
en production :

\begin{itemize}
  \item \textbf{Interface utilisateur} : tableau de bord (dashboard) pour les analystes
    permettant de consulter les alertes, de les valider ou de les rejeter, et de consulter
    les explications en langage naturel sans accéder au code.
  \item \textbf{Boucle de rétroaction (\textit{feedback loop})} : intégration des
    décisions humaines (fraude confirmée ou infirmée) pour ré-entraîner périodiquement
    les modèles et améliorer leur précision sur les données métier réelles.
  \item \textbf{Détection comportementale} : profiling par utilisateur ou par compte
    pour détecter non plus des transactions isolées anormales, mais des dérives
    comportementales dans le temps (ex : montants progressivement croissants, changements
    d'habitudes de transfert).
\end{itemize}

% ---------------------------------------------------
\section*{Conclusion}
\addcontentsline{toc}{section}{Conclusion}
% ---------------------------------------------------

Ce chapitre a présenté l'architecture de déploiement et les mécanismes de reproductibilité
du projet. La \textbf{structure modulaire} et l'orchestration via \texttt{run\_all.py}
fournissent une base maintenable et facilement reprise par un nouveau développeur. Les
quatre garanties de reproductibilité --- graines fixes, environnement versionné, artefacts
sérialisés, séparation étanche des données --- assurent que les résultats présentés dans
ce rapport sont intégralement reproductibles. La suite de tests structurée pose les bases
d'un contrôle qualité progressif. Les perspectives identifiées offrent de nombreuses pistes
pour prolonger ce travail vers un \textbf{système de production opérationnel}, capable
d'analyser des flux transactionnels en temps réel et d'apprendre de façon continue
des retours des auditeurs.


% =====================================================
% CONCLUSION GÉNÉRALE
% =====================================================
\chapter*{Conclusion générale}
\addcontentsline{toc}{chapter}{Conclusion générale}

Ce projet de fin d'études a abouti à un résultat concret : une \textbf{chaîne complète de
détection intelligente d'anomalies financières}, opérationnelle, reproductible et
expliquable, développée selon la méthodologie CRISP-DM sur des données transactionnelles
réelles issues d'une mission d'audit IT chez PwC Tunisie.

\paragraph{Récapitulatif de la démarche.} Partant d'un corpus de 200~000 transactions
financières, le travail a progressé de manière structurée : compréhension des enjeux
métier, analyse exploratoire approfondie, préparation rigoureuse des données (14 variables,
séparation étanche des partitions, transformation logarithmique, one-hot encoding),
entraînement de cinq configurations de modèles, évaluation comparative et mise en place
d'une double couche d'explicabilité (SHAP/LIME + Ollama).

\paragraph{Réponse à la problématique.} Le pipeline construit répond positivement à la
problématique initiale : il est \textbf{reproductible} (orchestration via \texttt{run\_all.py},
graines fixes, versions épinglées), \textbf{efficace} (la forêt aléatoire avec SMOTE
atteint un F1-score de $\sim$0,80 et un PR-AUC de $\sim$0,84 sur le test, contre
F1 = 0,0077 pour la baseline), et \textbf{interprétable} (chaque transaction suspecte est
accompagnée d'une explication en langage naturel via Ollama).

\paragraph{Difficultés et apports techniques.} Trois défis ont retenu l'essentiel de mon
attention : la gestion du déséquilibre extrême (ratio 1:774), la discipline anti-leakage
à maintenir tout au long du prétraitement, et la sensibilité du seuil de détection de
l'autoencoder. Sur le plan personnel, ce projet m'a permis de mettre en œuvre, sur un cas
réel et contraint, un spectre technique inhabituel pour un projet de fin d'études : machine
learning supervisé, deep learning, méthodes d'interprétabilité et intégration d'un LLM
local --- le tout dans un cadre méthodologique rigoureux imposé par le contexte professionnel.

\paragraph{Perspectives.} Ce travail ouvre plusieurs axes de prolongement : l'intégration
de modèles de gradient boosting (LightGBM, XGBoost), le passage à un autoencoder
variationnel, le déploiement en temps réel sur un flux de transactions, et l'enrichissement
de la couche d'explicabilité par des interfaces utilisateur dédiées aux analystes métier.

% =====================================================
% BIBLIOGRAPHIE
% =====================================================
\chapter*{Bibliographie / Netographie}
\addcontentsline{toc}{chapter}{Bibliographie / Netographie}

\begin{enumerate}

  \item \textbf{Lopez-Rojas, E.~A., Elmir, A., \& Axelsson, S.} (2016). PaySim: A
    financial mobile money simulator for fraud detection. \textit{The 28th European
    Modeling and Simulation Symposium (EMSS)}, Larnaca, Cyprus.

  \item \textbf{Chapman, P. et al.} (2000). CRISP-DM 1.0: Step-by-step data mining guide.
    SPSS Inc.

  \item \textbf{Chawla, N.~V. et al.} (2002). SMOTE: Synthetic minority over-sampling
    technique. \textit{Journal of Artificial Intelligence Research}, 16, 321--357.

  \item \textbf{Breiman, L.} (2001). Random forests. \textit{Machine Learning}, 45(1),
    5--32.

  \item \textbf{Hinton, G.~E., \& Salakhutdinov, R.~R.} (2006). Reducing the
    dimensionality of data with neural networks. \textit{Science}, 313(5786), 504--507.

  \item \textbf{Lundberg, S.~M., \& Lee, S.-I.} (2017). A unified approach to interpreting
    model predictions. \textit{NeurIPS}, 30.

  \item \textbf{Ribeiro, M.~T., Singh, S., \& Guestrin, C.} (2016). Why should I trust
    you? Explaining the predictions of any classifier. \textit{KDD 2016}, pp.~1135--1144.

  \item \textbf{Pedregosa, F. et al.} (2011). Scikit-learn: Machine Learning in Python.
    \textit{JMLR}, 12, 2825--2830. \url{https://scikit-learn.org/stable/}

  \item \textbf{Lemaître, G., Nogueira, F., \& Aridas, C.~K.} (2017). Imbalanced-learn.
    \textit{JMLR}, 18(17), 1--5. \url{https://imbalanced-learn.org/stable/}

  \item \textbf{Lundberg, S.~M.} Documentation officielle SHAP.
    \url{https://shap.readthedocs.io/}

  \item \textbf{Ribeiro, M.~T.} Documentation officielle LIME.
    \url{https://lime-ml.readthedocs.io/}

  \item \textbf{Ollama Inc.} Documentation officielle Ollama.
    \url{https://ollama.com/}

  \item \textbf{Chaieb, Y.} \textit{Intelligent Anomaly Detection} --- Dépôt GitHub
    du projet. PwC Tunisie / ESPRIT, 2026.
    \url{https://github.com/yousra0/Intelligent-Anomaly-Detection}

  \item \textbf{Auteur non spécifié.} La méthodologie CRISP-DM : présentation des six
    phases. Le1817.com. \url{https://www.le1817.com/blog/la-methodologie-crisp-dm}

\end{enumerate}

% =====================================================
% ANNEXES
% =====================================================
\appendix

\chapter{Hyperparamètres détaillés des modèles}

\section{Régression logistique}

\begin{lstlisting}[caption={Hyperparamètres LR (src/models/ml\_models.py)}]
LR_DEFAULTS = {
    "C":             0.1,
    "max_iter":      1000,
    "solver":        "lbfgs",
    "class_weight":  "balanced",
    "random_state":  42,
    "n_jobs":        -1,
}
\end{lstlisting}

\section{Forêt aléatoire}

\begin{lstlisting}[caption={Hyperparamètres RF (src/models/ml\_models.py)}]
RF_DEFAULTS = {
    "n_estimators":    300,
    "max_depth":       10,
    "min_samples_leaf": 5,
    "class_weight":    "balanced",
    "random_state":    42,
    "n_jobs":          -1,
}
\end{lstlisting}

\section{Autoencoder}

\begin{lstlisting}[caption={Hyperparamètres AE (src/models/autoencoder.py)}]
AE_DEFAULTS = {
    "encoder_dims":         [10, 7],
    "bottleneck_dim":       4,
    "decoder_dims":         [7, 10],
    "activation":           "relu",
    "output_activation":    "linear",
    "dropout_rate":         0.2,
    "use_batch_norm":       True,
    "l2_reg":               1e-5,
    "epochs":               100,
    "batch_size":           256,
    "learning_rate":        1e-3,
    "patience":             10,
    "val_split":            0.1,
    "threshold_percentile": 95,
}
\end{lstlisting}

\chapter{Pipeline de préparation --- variables}

\section{Variables normalisées (StandardScaler)}

\begin{lstlisting}
SCALE_COLS = [
    "step", "hour", "day", "week",
    "log_amount", "balance_diff_orig"
]
\end{lstlisting}

\section{Variables binaires (pas de normalisation)}

\begin{lstlisting}
BINARY_COLS = [
    "high_risk_hour",
    "is_transfer_or_cashout",
    "dest_zero_balance"
]
\end{lstlisting}

\section{Variables one-hot (type de transaction)}

\begin{lstlisting}
TYPE_COLS = [
    "type_CASH_IN", "type_CASH_OUT", "type_DEBIT",
    "type_PAYMENT", "type_TRANSFER"
]
\end{lstlisting}

\chapter{Planning de réalisation}

\begin{table}[H]
\centering
\caption{Planning de réalisation du projet de fin d'études}
\label{tab:planning}
\begin{tabular}{lll}
\toprule
\textbf{Phase CRISP-DM} & \textbf{Activités principales} & \textbf{Durée} \\
\midrule
Compréhension métier    & Analyse des besoins, définition des objectifs & Semaine 1 \\
Compréhension données   & EDA, notebook 01, analyse qualité & Semaines 2--3 \\
Préparation données     & Feature engineering, split, normalisation & Semaines 3--4 \\
Modélisation            & LR, RF, Autoencoder, réglage seuils & Semaines 5--7 \\
Évaluation              & Métriques, SHAP, LIME, Ollama & Semaines 7--9 \\
Déploiement             & Pipeline, tests, documentation & Semaines 9--10 \\
Rédaction du rapport    & Rédaction, révision, mise en forme & Semaines 8--12 \\
\bottomrule
\end{tabular}
\end{table}

\chapter{Exemple d'explication générée par Ollama}

\begin{lstlisting}[caption={Exemple d'explication pour une transaction suspecte}]
Transaction ID : TXN_00123
Montant : 245 000 | Type : TRANSFER | Heure : 03h00

Cette transaction est signalee comme suspecte pour les raisons suivantes :
1. Type TRANSFER : seul type associe a des fraudes dans le corpus d'audit.
2. Heure 3h00 : heure a risque eleve (ratio fraude 10x superieur a la
   moyenne horaire).
3. Difference de solde emetteur (balance_diff_orig = 245 000) anormalement
   elevee, coherente avec un virement frauduleux.
4. Score de reconstruction AE (MSE = 1.87) superieur au seuil (theta = 1.67)
   : pattern inhabituel non reconnu par le modele.

Recommandation : Soumettre a une verification manuelle par un auditeur.
\end{lstlisting}

\end{document}
